\documentclass[11pt]{article}

% ---- theme variables (passed via pandoc -V) ----
\newcommand{\ThemeName}{Git \& Dev Workflow}
\newcommand{\ThemeKicker}{Quality Track}
\newcommand{\ThemeNumber}{}

\usepackage{interviewstyle}
\setthemecolor{ea580c}{fb923c}

% pandoc-required plumbing
\providecommand{\tightlist}{\setlength{\itemsep}{1pt}\setlength{\parskip}{0pt}}
\setlength{\emergencystretch}{3em}
\providecommand{\pandocbounded}[1]{#1}

% syntax-highlighting macros emitted by pandoc (skylighting)
\usepackage{color}
\usepackage{fancyvrb}
\newcommand{\VerbBar}{|}
\newcommand{\VERB}{\Verb[commandchars=\\\{\}]}
\DefineVerbatimEnvironment{Highlighting}{Verbatim}{commandchars=\\\{\}}
% Add ',fontsize=\small' for more characters per line
\usepackage{framed}
\definecolor{shadecolor}{RGB}{35,38,41}
\newenvironment{Shaded}{\begin{snugshade}}{\end{snugshade}}
\newcommand{\AlertTok}[1]{\textcolor[rgb]{0.58,0.85,0.30}{\textbf{\colorbox[rgb]{0.30,0.12,0.14}{#1}}}}
\newcommand{\AnnotationTok}[1]{\textcolor[rgb]{0.25,0.50,0.35}{#1}}
\newcommand{\AttributeTok}[1]{\textcolor[rgb]{0.16,0.50,0.73}{#1}}
\newcommand{\BaseNTok}[1]{\textcolor[rgb]{0.96,0.45,0.00}{#1}}
\newcommand{\BuiltInTok}[1]{\textcolor[rgb]{0.50,0.55,0.55}{#1}}
\newcommand{\CharTok}[1]{\textcolor[rgb]{0.24,0.68,0.91}{#1}}
\newcommand{\CommentTok}[1]{\textcolor[rgb]{0.48,0.49,0.49}{#1}}
\newcommand{\CommentVarTok}[1]{\textcolor[rgb]{0.50,0.55,0.55}{#1}}
\newcommand{\ConstantTok}[1]{\textcolor[rgb]{0.15,0.68,0.68}{\textbf{#1}}}
\newcommand{\ControlFlowTok}[1]{\textcolor[rgb]{0.99,0.74,0.29}{\textbf{#1}}}
\newcommand{\DataTypeTok}[1]{\textcolor[rgb]{0.16,0.50,0.73}{#1}}
\newcommand{\DecValTok}[1]{\textcolor[rgb]{0.96,0.45,0.00}{#1}}
\newcommand{\DocumentationTok}[1]{\textcolor[rgb]{0.64,0.20,0.25}{#1}}
\newcommand{\ErrorTok}[1]{\textcolor[rgb]{0.85,0.27,0.33}{\underline{#1}}}
\newcommand{\ExtensionTok}[1]{\textcolor[rgb]{0.00,0.60,1.00}{\textbf{#1}}}
\newcommand{\FloatTok}[1]{\textcolor[rgb]{0.96,0.45,0.00}{#1}}
\newcommand{\FunctionTok}[1]{\textcolor[rgb]{0.56,0.27,0.68}{#1}}
\newcommand{\ImportTok}[1]{\textcolor[rgb]{0.15,0.68,0.38}{#1}}
\newcommand{\InformationTok}[1]{\textcolor[rgb]{0.77,0.36,0.00}{#1}}
\newcommand{\KeywordTok}[1]{\textcolor[rgb]{0.81,0.81,0.76}{\textbf{#1}}}
\newcommand{\NormalTok}[1]{\textcolor[rgb]{0.81,0.81,0.76}{#1}}
\newcommand{\OperatorTok}[1]{\textcolor[rgb]{0.81,0.81,0.76}{#1}}
\newcommand{\OtherTok}[1]{\textcolor[rgb]{0.15,0.68,0.38}{#1}}
\newcommand{\PreprocessorTok}[1]{\textcolor[rgb]{0.15,0.68,0.38}{#1}}
\newcommand{\RegionMarkerTok}[1]{\textcolor[rgb]{0.16,0.50,0.73}{\colorbox[rgb]{0.08,0.19,0.26}{#1}}}
\newcommand{\SpecialCharTok}[1]{\textcolor[rgb]{0.24,0.68,0.91}{#1}}
\newcommand{\SpecialStringTok}[1]{\textcolor[rgb]{0.85,0.27,0.33}{#1}}
\newcommand{\StringTok}[1]{\textcolor[rgb]{0.96,0.31,0.31}{#1}}
\newcommand{\VariableTok}[1]{\textcolor[rgb]{0.15,0.68,0.68}{#1}}
\newcommand{\VerbatimStringTok}[1]{\textcolor[rgb]{0.85,0.27,0.33}{#1}}
\newcommand{\WarningTok}[1]{\textcolor[rgb]{0.85,0.27,0.33}{#1}}

\begin{document}

\makecoverpage

\thispagestyle{empty}
{\hypersetup{hidelinks}\tableofcontents}
\clearpage

\begin{quote}
\textbf{How to use this file:} Read top-to-bottom for deep, mechanism-level understanding. Jump to §Common Interview Questions and §Revision Cheat Sheet for last-minute revision. Git fluency is \emph{assumed} on the job --- the differentiator in interviews is that you understand the \textbf{object model} underneath the commands, can reason about what \texttt{reset\ -\/-hard} actually touches, and can recover from disasters calmly. Memorizing flags is worthless; understanding that "a branch is just a 41-byte file containing a hash" is everything.
\end{quote}

\section{Overview --- What It Is}\label{overview--what-it-is}

\textbf{Git} is a \textbf{distributed version control system (DVCS)}. Every clone is a \emph{complete, independent copy} of the entire repository --- all history, all branches, all tags --- not just the latest snapshot. There is no privileged central server in Git\textquotesingle s design; "origin" is just a convention. You can commit, branch, merge, view history, and time-travel entirely offline.

At its heart, Git is shockingly simple. Linus Torvalds famously described it as a \textbf{"content-addressable filesystem with a version-control UI bolted on top."} Strip away the porcelain commands (\texttt{add}, \texttt{commit}, \texttt{merge}) and you find a tiny \textbf{key-value store}: you give Git some content, it gives you back a hash (the key); you give it the hash later, it gives you back the exact content. Everything else --- commits, branches, tags, history --- is built from that one primitive.

\setcodelang{}

\begin{verbatim}
working dir  ──git add──▶  staging (index)  ──git commit──▶  local repo  ──push──▶  remote
  (edit)                    (stage)                          (.git/ history)         (share)
   ▲                                                              │
   └──────────────────── git checkout / restore ◀────────────────┘
\end{verbatim}

Three things make Git fundamentally different from what came before:

\begin{enumerate}
\def\labelenumi{\arabic{enumi}.}
\tightlist
\item
  \textbf{Snapshots, not diffs.} Most older systems (SVN, CVS, Perforce) store a file and then a \emph{list of changes} to it over time. Git stores a \textbf{full snapshot} of the whole project at each commit (with deduplication so it\textquotesingle s not wasteful). This makes branching and merging cheap and makes operations like \texttt{checkout} and \texttt{diff} between arbitrary points fast.
\item
  \textbf{Content addressing.} Every object (file content, directory tree, commit) is named by the \textbf{SHA-1 hash of its content}. Identical content \(\rightarrow\) identical hash \(\rightarrow\) stored once. This gives you free deduplication and tamper-evidence.
\item
  \textbf{Distributed.} Every clone is a backup. There is no single point of failure for \emph{history}. The "server" (GitHub/GitLab) adds collaboration features (PRs, access control, CI) but Git itself doesn\textquotesingle t require one.
\end{enumerate}

\setcodelang{}

\begin{verbatim}
       Centralized (SVN)                    Distributed (Git)
   ┌──────────────────────┐          ┌──────────────────────────┐
   │   central server     │          │  origin (a clone, by      │
   │  (the ONLY history)  │          │  convention "the truth")  │
   └─────────┬────────────┘          └───────┬──────────────────┘
             │ checkout (files only)         │ clone (FULL history)
     ┌───────┼───────┐                ┌───────┼────────┐
     ▼       ▼       ▼                ▼        ▼        ▼
   dev1    dev2    dev3            dev1      dev2     dev3
 (no local history)            (each has FULL history; can
                                commit/branch/merge offline)
\end{verbatim}

\section{Why It Exists}\label{why-it-exists}

Coordinating many people editing the same codebase --- without overwriting each other, while still being able to \textbf{undo, branch, review, audit, and bisect} --- is impossible to do by hand (copying \texttt{project\_final\_FINAL\_v2/} folders does not scale and has no merge story).

\textbf{The problems version control solves:}

\begin{enumerate}
\def\labelenumi{\arabic{enumi}.}
\tightlist
\item
  \textbf{Concurrent editing.} Two people change the same file. Without VCS, one overwrites the other. Git lets both work in isolation and \emph{merges} the changes, only stopping to ask a human when the edits truly conflict (touch the same lines).
\item
  \textbf{History \& accountability.} Who changed this line, when, and \emph{why}? \texttt{git\ blame} + the commit message answers this. Critical for debugging ("this broke in the deploy three weeks ago") and for institutional memory.
\item
  \textbf{Safe experimentation.} Cheap branches mean you can try a risky refactor, and if it fails, delete the branch --- main is untouched.
\item
  \textbf{Reproducibility.} A commit hash pins the \emph{exact} state of the codebase. CI builds, deploys, and rollbacks all reference immutable commits.
\item
  \textbf{Distributed resilience.} Every developer\textquotesingle s laptop is a full backup. GitHub going down doesn\textquotesingle t lose your history.
\end{enumerate}

\textbf{Why Git specifically (vs. earlier VCS):}

\begin{longtable}[]{@{}
  >{\raggedright\arraybackslash}p{(\columnwidth - 4\tabcolsep) * \real{0.3026}}
  >{\raggedright\arraybackslash}p{(\columnwidth - 4\tabcolsep) * \real{0.3684}}
  >{\raggedright\arraybackslash}p{(\columnwidth - 4\tabcolsep) * \real{0.3289}}@{}}
\toprule\noalign{}
\rowcolor{acc!16}
\begin{minipage}[b]{\linewidth}\raggedright
Need
\end{minipage} & \begin{minipage}[b]{\linewidth}\raggedright
SVN/CVS (centralized)
\end{minipage} & \begin{minipage}[b]{\linewidth}\raggedright
Git (distributed)
\end{minipage} \\
\midrule\noalign{}
\endhead
\bottomrule\noalign{}
\endlastfoot
Commit offline & No (needs server) & Yes \\
\rowcolor{zebra}
Branch creation cost & Expensive (server-side copy) & \textasciitilde Instant (write one file) \\
Merge quality & Painful, weak & Strong 3-way merge \\
\rowcolor{zebra}
Full history locally & No & Yes \\
Tamper-evidence & No & Yes (hash-chained) \\
\rowcolor{zebra}
Single point of failure & Yes (the server) & No (for history) \\
\end{longtable}

\begin{quote}
Git stores \textbf{content addressed by hash}, so identical content is stored exactly once and history is tamper-evident: change any old commit and \emph{every later hash changes}, instantly visible to everyone.
\end{quote}

\section{Why FAANG Cares}\label{why-faang-cares}

Git is the daily substrate of every engineer\textquotesingle s work. Interviews rarely ask "what does \texttt{git\ add} do" --- they probe whether you understand the \emph{model} well enough to (a) keep history clean for reviewers, (b) recover from mistakes without panicking, and (c) reason about workflow at scale.

\textbf{Google:}

\begin{itemize}
\tightlist
\item
  Runs one of the largest \textbf{monorepos} on earth (\texttt{google3}), backed by Piper (a Perforce-derived system) and the Mondrian/Critique review tooling --- but external repos and Android (AOSP, which uses Git + \texttt{repo}) are pure Git. Engineers are expected to produce small, reviewable, well-described CLs (changelists). The cultural value: \emph{clean, atomic, bisectable history}. Knowing how \texttt{git\ bisect} finds a regression in O(log n) is directly relevant to debugging at Google scale.
\end{itemize}

\textbf{Meta:}

\begin{itemize}
\tightlist
\item
  Also a giant \textbf{monorepo} (backed by Mercurial-derived \textbf{Sapling/EdenFS} with \texttt{sl}, conceptually similar to Git). Trunk-based development with thousands of commits/day. They invest heavily in \emph{virtual filesystems} and \emph{sparse checkouts} because the repo is too big to fully materialize. Phabricator/Diffstack-style stacked diffs are core. Understanding why a monorepo needs sparse-checkout and VFS shows systems maturity.
\end{itemize}

\textbf{Amazon:}

\begin{itemize}
\tightlist
\item
  Heavy \textbf{polyrepo} culture (thousands of small service repos), tied to internal CI/CD (Pipelines, Brazil build system, CodeCommit). Two-pizza teams own their repos end-to-end. Interviews and on-the-job reviews value clean PRs, conventional commit hygiene, and the ability to cherry-pick hotfixes onto release branches.
\end{itemize}

\textbf{Apple:}

\begin{itemize}
\tightlist
\item
  Mixed; strong emphasis on commit signing, provenance, and careful release branching for OS versions. GPG/SSH-signed commits and tags matter for supply-chain trust.
\end{itemize}

\textbf{Microsoft (GitHub, Azure DevOps):}

\begin{itemize}
\tightlist
\item
  \emph{Owns} GitHub. Migrated Windows (\textasciitilde300 GB, 3.5M files) to Git, which required inventing \textbf{VFS for Git} (now \textbf{Scalar}) and partial clone. The Windows-on-Git story is a famous case study in scaling Git. Azure DevOps and GitHub Actions are their CI/CD.
\end{itemize}

\textbf{Netflix / Uber / Stripe / Databricks:}

\begin{itemize}
\tightlist
\item
  Trunk-based development, feature flags, heavy CI gates, and "you build it, you run it." Clean Git workflow is part of operational safety --- a botched force-push can take down a deploy pipeline.
\end{itemize}

\textbf{The interview reality:} When you say "I\textquotesingle d cherry-pick the fix onto the release branch," or "I\textquotesingle ll rebase my feature branch to keep history linear, but I\textquotesingle ll never rebase a shared branch," you\textquotesingle re signaling that you understand collaboration safety and history hygiene --- the things that separate a senior from a junior who only knows \texttt{git\ add\ .\ \&\&\ git\ commit\ -m\ "stuff"}.

\section{Core Concepts}\label{core-concepts}

\subsection{Snapshots, Not Diffs (the mental-model shift)}\label{snapshots-not-diffs-the-mental-model-shift}

This is the single most important conceptual leap. Imagine three commits where only one file changes each time:

\setcodelang{}

\begin{verbatim}
SVN/CVS style (store deltas):                Git style (store snapshots):
  C1: full A, full B, full C                   C1: ┌snap┐ A1  B1  C1
  C2: Δ(A only)                                C2: ┌snap┐ A2  B1* C1*   (* = reuse, same hash)
  C3: Δ(B only)                                C3: ┌snap┐ A2* B2  C1*

  To reconstruct C3 you must               To reconstruct C3, just read its
  replay C1 + all deltas.                  snapshot. Unchanged files point to
                                           the SAME blob object (dedup).
\end{verbatim}

In Git, a commit references a complete \textbf{tree} (the whole directory structure). Files that didn\textquotesingle t change between commits \textbf{point to the exact same blob object} --- same content, same hash, stored once. So you get the conceptual simplicity of "every commit is a full snapshot" with the storage efficiency of deltas (added later via packfiles). This is why \texttt{git\ checkout\ \textless{}old-commit\textgreater{}} is fast and why \texttt{git\ diff} between any two arbitrary commits is cheap.

\subsection{Content Addressing}\label{content-addressing}

Every piece of data Git stores is named by the \textbf{hash of its content}, not by a filename or a sequential ID. The address \emph{is} derived from the bytes. Consequences:

\begin{itemize}
\tightlist
\item
  \textbf{Deduplication is automatic.} Two files with identical bytes anywhere in history \(\rightarrow\) one stored object.
\item
  \textbf{Integrity is automatic.} If a byte on disk corrupts, the content no longer hashes to its name --- Git detects it (\texttt{git\ fsck}).
\item
  \textbf{History is tamper-evident.} Because a commit\textquotesingle s hash includes its parent\textquotesingle s hash (a hash chain, like a blockchain), altering any historical commit changes its hash, which changes its child\textquotesingle s hash, cascading to every descendant. You cannot quietly rewrite the past.
\end{itemize}

\subsection{A Branch Is Just a Pointer}\label{a-branch-is-just-a-pointer}

The most liberating realization: \textbf{a branch is a 41-byte text file} (40 hex chars + newline) at \texttt{.git/refs/heads/\textless{}name\textgreater{}} containing one commit hash. Creating a branch writes one tiny file. Switching branches just moves \texttt{HEAD}. That\textquotesingle s why branching in Git is instantaneous, unlike SVN where a branch was a server-side directory copy.

\setcodelang{Bash}

\begin{Shaded}
\begin{Highlighting}[]
\ExtensionTok{$}\NormalTok{ cat .git/refs/heads/main}
\ExtensionTok{3f8a1c9e2b...d4}    \CommentTok{\# that\textquotesingle{}s the entire branch. one hash.}
\ExtensionTok{$}\NormalTok{ cat .git/HEAD}
\ExtensionTok{ref:}\NormalTok{ refs/heads/main   }\CommentTok{\# HEAD points to the current branch}
\end{Highlighting}
\end{Shaded}

\subsection{Commits Are Immutable; "Editing History" Means Rewriting}\label{commits-are-immutable-editing-history-means-rewriting}

You never \emph{modify} a commit. Operations that appear to edit history (\texttt{amend}, \texttt{rebase}, \texttt{commit\ -\/-fixup}) actually \textbf{create brand-new commits with new hashes} and move a branch pointer to them. The old commits become unreferenced (dangling) and are eventually garbage-collected. This is why rebasing shared history is dangerous (covered later) --- the old hashes others depend on vanish.

\section{The Git Object Model}\label{the-git-object-model}

Everything in \texttt{.git/objects/} is one of \textbf{four object types}. Understanding these four is understanding Git.

\begin{longtable}[]{@{}
  >{\raggedright\arraybackslash}p{(\columnwidth - 6\tabcolsep) * \real{0.1357}}
  >{\raggedright\arraybackslash}p{(\columnwidth - 6\tabcolsep) * \real{0.4392}}
  >{\raggedright\arraybackslash}p{(\columnwidth - 6\tabcolsep) * \real{0.1809}}
  >{\raggedright\arraybackslash}p{(\columnwidth - 6\tabcolsep) * \real{0.2442}}@{}}
\toprule\noalign{}
\rowcolor{acc!16}
\begin{minipage}[b]{\linewidth}\raggedright
Object
\end{minipage} & \begin{minipage}[b]{\linewidth}\raggedright
Stores
\end{minipage} & \begin{minipage}[b]{\linewidth}\raggedright
Named by
\end{minipage} & \begin{minipage}[b]{\linewidth}\raggedright
Points to
\end{minipage} \\
\midrule\noalign{}
\endhead
\bottomrule\noalign{}
\endlastfoot
\textbf{blob} & file \emph{contents} (no name, no mode) & SHA-1 of content & nothing \\
\rowcolor{zebra}
\textbf{tree} & a directory listing (names, modes, \(\rightarrow\) blobs/trees) & SHA-1 of its entries & blobs and subtrees \\
\textbf{commit} & snapshot pointer + metadata & SHA-1 of commit text & one tree + parent commit(s) \\
\rowcolor{zebra}
\textbf{tag} (annotated) & a named, possibly signed pointer to an object & SHA-1 of tag text & usually a commit \\
\end{longtable}

\setcodelang{}

\begin{verbatim}
                    commit  3f8a1c…
                    ┌─────────────────────────────┐
                    │ tree    a1b2c3…              │───────┐
                    │ parent  9e8d7c… (prev commit)│──┐    │
                    │ author  Asha <a@x> 1718…     │  │    ▼
                    │ message "Add login handler"  │  │  tree a1b2c3…  (root dir)
                    └─────────────────────────────┘  │  ┌──────────────────────────┐
                                                      │  │ 100644 blob d4e5… README │──▶ blob (file bytes)
                            (parent commit)◀──────────┘  │ 040000 tree f6a7… src/   │──┐
                                                         └──────────────────────────┘  │
                                                                                        ▼
                                                                            tree f6a7…  (src/ dir)
                                                                            ┌────────────────────────┐
                                                                            │ 100644 blob 0c1d… app.py│──▶ blob
                                                                            └────────────────────────┘
\end{verbatim}

\begin{itemize}
\tightlist
\item
  \textbf{blob} = the contents of a file. Note: a blob has \emph{no filename}. The filename lives in the \textbf{tree} that references it. So renaming a file (same content) doesn\textquotesingle t create a new blob --- it changes a tree entry.
\item
  \textbf{tree} = a directory. It maps \texttt{mode\ +\ name\ →\ blob/tree\ hash}. The root tree of a commit captures the whole project layout.
\item
  \textbf{commit} = a tree (the snapshot) + zero-or-more \textbf{parents} + author/committer + timestamp + message. The first commit has zero parents; a normal commit has one; a merge commit has two or more.
\item
  \textbf{tag} = a permanent named pointer. \emph{Lightweight} tags are just refs; \emph{annotated} tags are real objects (taggable, signable, with their own message).
\end{itemize}

\subsection{How a Commit Hash Is Computed (worked)}\label{how-a-commit-hash-is-computed-worked}

A Git object\textquotesingle s hash is the SHA-1 of: \texttt{"\textless{}type\textgreater{}\ \textless{}byte-length\textgreater{}\textbackslash{}0"\ +\ content}. Let\textquotesingle s compute a blob hash by hand-equivalent:

\setcodelang{Bash}

\begin{Shaded}
\begin{Highlighting}[]
\ExtensionTok{$}\NormalTok{ printf }\StringTok{\textquotesingle{}hello\textbackslash{}n\textquotesingle{}} \KeywordTok{|} \FunctionTok{git}\NormalTok{ hash{-}object }\AttributeTok{{-}{-}stdin}
\ExtensionTok{ce013625030ba8dba906f756967f9e9ca394464a}

\CommentTok{\# Reproduce it manually: the header is "blob 6\textbackslash{}0" (6 = len of "hello\textbackslash{}n")}
\ExtensionTok{$}\NormalTok{ printf }\StringTok{\textquotesingle{}blob 6\textbackslash{}0hello\textbackslash{}n\textquotesingle{}} \KeywordTok{|} \FunctionTok{sha1sum}
\ExtensionTok{ce013625030ba8dba906f756967f9e9ca394464a}   \AttributeTok{{-}}    \CommentTok{\# identical!}
\end{Highlighting}
\end{Shaded}

A \textbf{commit} hash is the SHA-1 of the commit\textquotesingle s text, which literally includes the tree hash, the parent hash, author/committer lines, and the message:

\setcodelang{Bash}

\begin{Shaded}
\begin{Highlighting}[]
\ExtensionTok{$}\NormalTok{ git cat{-}file }\AttributeTok{{-}p}\NormalTok{ HEAD}
\ExtensionTok{tree}\NormalTok{ a1b2c3d4e5f6a7b8c9d0e1f2a3b4c5d6e7f8a9b0}
\ExtensionTok{parent}\NormalTok{ 9e8d7c6b5a4f3e2d1c0b9a8f7e6d5c4b3a2f1e0d}
\ExtensionTok{author}\NormalTok{  Asha Rao }\OperatorTok{\textless{}}\NormalTok{asha@x.com}\OperatorTok{\textgreater{}}\NormalTok{ 1718000000 +0000}
\ExtensionTok{committer}\NormalTok{ Asha Rao }\OperatorTok{\textless{}}\NormalTok{asha@x.com}\OperatorTok{\textgreater{}}\NormalTok{ 1718000000 +0000}

\ExtensionTok{Add}\NormalTok{ login handler}
\end{Highlighting}
\end{Shaded}

Because the \textbf{parent hash is part of the commit text}, changing any ancestor changes that commit\textquotesingle s text \(\rightarrow\) its hash \(\rightarrow\) and therefore the hashes of \emph{every} descendant. That cascade is the tamper-evidence property. It also means: \textbf{the same diff applied at a different point in history produces a different commit hash} (the parent differs). That single fact explains everything about rebase, cherry-pick, and "why rebasing shared history breaks people."

\subsection{Identical Content Dedupes --- demonstrated}\label{identical-content-dedupes--demonstrated}

\setcodelang{Bash}

\begin{Shaded}
\begin{Highlighting}[]
\ExtensionTok{$}\NormalTok{ echo }\StringTok{"config"} \OperatorTok{\textgreater{}}\NormalTok{ a.txt}
\ExtensionTok{$}\NormalTok{ echo }\StringTok{"config"} \OperatorTok{\textgreater{}}\NormalTok{ b.txt}
\ExtensionTok{$}\NormalTok{ git add a.txt b.txt}
\ExtensionTok{$}\NormalTok{ git ls{-}files }\AttributeTok{{-}s}
\ExtensionTok{100644}\NormalTok{ 5e40c0877058c504203932e5136051cf3cd3519b 0   a.txt}
\ExtensionTok{100644}\NormalTok{ 5e40c0877058c504203932e5136051cf3cd3519b 0   b.txt}
\CommentTok{\#              \^{}\^{}\^{}\^{} SAME hash — both filenames point to ONE blob object.}
\end{Highlighting}
\end{Shaded}

Two files, identical content \(\rightarrow\) one blob stored. The tree records two \emph{names} pointing at the same blob. This is automatic; you never ask for it.

\subsection{\texorpdfstring{\texttt{git\ cat-file} \& \texttt{git\ hash-object} --- looking inside}{git cat-file \& git hash-object --- looking inside}}\label{git-cat-file--git-hash-object--looking-inside}

\setcodelang{Bash}

\begin{Shaded}
\begin{Highlighting}[]
\ExtensionTok{$}\NormalTok{ git cat{-}file }\AttributeTok{{-}t}\NormalTok{ 3f8a1c   }\CommentTok{\# type of an object}
\ExtensionTok{commit}
\ExtensionTok{$}\NormalTok{ git cat{-}file }\AttributeTok{{-}s}\NormalTok{ 3f8a1c   }\CommentTok{\# size in bytes}
\ExtensionTok{176}
\ExtensionTok{$}\NormalTok{ git cat{-}file }\AttributeTok{{-}p}\NormalTok{ a1b2c3   }\CommentTok{\# pretty{-}print (here, a tree)}
\ExtensionTok{100644}\NormalTok{ blob d4e5f6…  README.md}
\ExtensionTok{040000}\NormalTok{ tree f6a7b8…  src}

\CommentTok{\# Store arbitrary content and get its hash (without committing):}
\ExtensionTok{$}\NormalTok{ echo }\StringTok{"test"} \KeywordTok{|} \FunctionTok{git}\NormalTok{ hash{-}object }\AttributeTok{{-}w} \AttributeTok{{-}{-}stdin}
\ExtensionTok{9daeafb9864cf43055ae93beb0afd6c7d144bfa4}   \CommentTok{\# {-}w actually writes it to .git/objects}
\end{Highlighting}
\end{Shaded}

\subsection{Packfiles \& Delta Compression}\label{packfiles--delta-compression}

Storing every snapshot as a full zlib-compressed blob ("loose objects") is fine for a while, but Git periodically runs \texttt{git\ gc}, which packs loose objects into a \textbf{packfile} (\texttt{*.pack} + \texttt{*.idx}). \emph{Within} a packfile, Git uses \textbf{delta compression}: similar objects (e.g., consecutive versions of a big file) are stored as one base + a chain of deltas. So Git\textquotesingle s user-facing model is "snapshots," but its \emph{storage} layer opportunistically uses deltas --- best of both worlds. Note: Git chooses deltas by \emph{content similarity heuristics} (size, filename), not by commit lineage; a delta base can even be a \emph{newer} object.

\setcodelang{Bash}

\begin{Shaded}
\begin{Highlighting}[]
\ExtensionTok{$}\NormalTok{ git gc                 }\CommentTok{\# pack loose objects}
\ExtensionTok{$}\NormalTok{ git count{-}objects }\AttributeTok{{-}vH}  \CommentTok{\# see loose vs packed counts and sizes}
\ExtensionTok{$}\NormalTok{ git verify{-}pack }\AttributeTok{{-}v}\NormalTok{ .git/objects/pack/pack{-}}\PreprocessorTok{*}\NormalTok{.idx }\KeywordTok{|} \FunctionTok{head}   \CommentTok{\# see delta chains}
\end{Highlighting}
\end{Shaded}

\subsection{Refs and HEAD}\label{refs-and-head}

A \textbf{ref} is a human name \(\rightarrow\) hash mapping, stored as a tiny file under \texttt{.git/refs/}:

\setcodelang{}

\begin{verbatim}
.git/
├── HEAD                      → "ref: refs/heads/main"  (the current branch)
├── refs/
│   ├── heads/                local branches
│   │   ├── main              → 3f8a1c…
│   │   └── feature/login     → 7c4b2a…
│   ├── tags/                 tags
│   │   └── v1.2.0            → e2f3a4…
│   └── remotes/origin/       remote-tracking branches (read-only mirrors)
│       └── main              → 3f8a1c…
\end{verbatim}

\begin{itemize}
\tightlist
\item
  \textbf{HEAD} is "where am I?" Usually it\textquotesingle s a \emph{symbolic ref} pointing at a branch (\texttt{ref:\ refs/heads/main}). When you commit, Git appends a commit and moves the branch HEAD points to.
\item
  \textbf{Detached HEAD} = HEAD points \emph{directly} at a commit hash instead of a branch. New commits you make have no branch tracking them --- easy to lose. (Recovery covered in Disasters.)
\item
  \textbf{Remote-tracking refs} (\texttt{origin/main}) are local read-only snapshots of where the remote was at your last \texttt{fetch}. \texttt{git\ fetch} updates them; \texttt{git\ push} updates the remote and then your tracking ref.
\end{itemize}

\begin{quote}
\textbf{Interview takeaway:} "Git is a content-addressed key-value store of four object types (blob, tree, commit, tag); branches and HEAD are just movable pointers (refs) into that graph." If you can say that sentence and explain the commit-hash-includes-parent cascade, you\textquotesingle ve demonstrated you understand Git rather than memorized it.
\end{quote}

\subsection{\texorpdfstring{SHA-1 \(\rightarrow\) SHA-256}{SHA-1 \textbackslash rightarrow SHA-256}}\label{sha-1-rightarrow-sha-256}

Git originally used \textbf{SHA-1} (160-bit). After the 2017 SHAttered SHA-1 collision, Git added two defenses: (1) a \textbf{collision-detection} variant of SHA-1 (\texttt{sha1dc}) that aborts on known-collision patterns, and (2) optional \textbf{SHA-256} repositories (\texttt{git\ init\ -\/-object-format=sha256}). SHA-256 isn\textquotesingle t yet the default because the ecosystem (hosting, tooling) is still migrating, and interoperability between SHA-1 and SHA-256 repos is a hard transition problem. For interviews: know that the hash is a content integrity mechanism, not a security boundary against a determined attacker --- that\textquotesingle s what \emph{signing} is for.

\section{The Three Areas \& The Index}\label{the-three-areas--the-index}

Git mediates between three "places," and the \textbf{index} (a.k.a. \textbf{staging area}) is the often-misunderstood middle layer that makes Git\textquotesingle s commit model precise.

\setcodelang{}

\begin{verbatim}
┌───────────────┐   git add        ┌───────────────┐   git commit    ┌───────────────┐
│  Working Dir  │ ───────────────▶ │  Index/Stage  │ ──────────────▶ │  Repository   │
│ (your files,  │                  │ (.git/index — │                 │ (.git/objects │
│  edited)      │ ◀─────────────── │  next snapshot)│ ◀───────────── │  — commits)   │
└───────────────┘  git restore     └───────────────┘   reset --mixed └───────────────┘
        ▲                                                                     │
        └──────────────── git checkout <commit> / restore --source ──────────┘
\end{verbatim}

\begin{itemize}
\tightlist
\item
  \textbf{Working directory} --- the actual files on disk you edit. This is the only place you see "uncommitted changes."
\item
  \textbf{Index (staging area)} --- a \emph{binary file} \texttt{.git/index} that holds the \textbf{proposed next commit}: a flat list of \texttt{(mode,\ blob-hash,\ stage,\ path)} entries. It is, essentially, a pre-baked tree. \texttt{git\ add} writes the file\textquotesingle s blob into the object store \emph{and} records its hash in the index.
\item
  \textbf{Repository} --- \texttt{.git/objects/}, the permanent immutable history.
\end{itemize}

The index is what enables \textbf{partial / staged commits}: you can edit five files but stage only two, producing a focused commit. The classic confusing diagram from \texttt{git\ status} ("Changes to be committed" vs "Changes not staged for commit") is literally describing index-vs-working-dir differences.

\subsection{Moving things between the areas}\label{moving-things-between-the-areas}

\setcodelang{Bash}

\begin{Shaded}
\begin{Highlighting}[]
\FunctionTok{git}\NormalTok{ add file.py              }\CommentTok{\# working dir → index (stage)}
\FunctionTok{git}\NormalTok{ add }\AttributeTok{{-}p}                   \CommentTok{\# interactively stage SOME hunks of a file (see below)}
\FunctionTok{git}\NormalTok{ restore }\AttributeTok{{-}{-}staged}\NormalTok{ file.py }\CommentTok{\# index → working dir  (UNSTAGE; keeps your edits)}
\FunctionTok{git}\NormalTok{ restore file.py          }\CommentTok{\# repo (HEAD) → working dir (DISCARD edits — destructive!)}
\FunctionTok{git}\NormalTok{ restore }\AttributeTok{{-}{-}source}\OperatorTok{=}\NormalTok{HEAD\textasciitilde{}2 file.py   }\CommentTok{\# pull an old version of a file into working dir}
\FunctionTok{git}\NormalTok{ rm }\AttributeTok{{-}{-}cached}\NormalTok{ file.py      }\CommentTok{\# stop tracking but keep the file on disk}
\FunctionTok{git}\NormalTok{ reset HEAD file.py       }\CommentTok{\# older syntax for "unstage" (= restore {-}{-}staged)}
\end{Highlighting}
\end{Shaded}

Modern Git (2.23+) split the overloaded \texttt{git\ checkout} into two clearer verbs:

\begin{itemize}
\tightlist
\item
  \textbf{\texttt{git\ switch}} --- change branches.
\item
  \textbf{\texttt{git\ restore}} --- restore file contents (from index or a commit).
\end{itemize}

\texttt{checkout} still works for both but is ambiguous (it does branches \emph{and} files), which historically caused accidental data loss when people typed \texttt{git\ checkout\ file} meaning to discard changes.

\subsection{\texorpdfstring{Partial staging with \texttt{git\ add\ -p}}{Partial staging with git add -p}}\label{partial-staging-with-git-add--p}

\texttt{add\ -p} (patch mode) walks you through each \emph{hunk} and asks whether to stage it:

\setcodelang{Bash}

\begin{Shaded}
\begin{Highlighting}[]
\ExtensionTok{$}\NormalTok{ git add }\AttributeTok{{-}p}\NormalTok{ app.py}
\ExtensionTok{@@} \AttributeTok{{-}10,7}\NormalTok{ +10,9 @@ def handler}\ErrorTok{(}\KeywordTok{)}\BuiltInTok{:}
\ExtensionTok{{-}}\NormalTok{    return None}
\ExtensionTok{+}\NormalTok{    log.info}\ErrorTok{(}\StringTok{"handling"}\KeywordTok{)}      \CommentTok{\# the bug fix you want NOW}
\ExtensionTok{+}\NormalTok{    return process}\ErrorTok{(}\ExtensionTok{req}\KeywordTok{)}
\ExtensionTok{@@} \AttributeTok{{-}40,3}\NormalTok{ +42,6 @@ def debug}\ErrorTok{(}\KeywordTok{)}\BuiltInTok{:}
\ExtensionTok{+}\NormalTok{    breakpoint}\ErrorTok{(}\KeywordTok{)}              \CommentTok{\# a stray debug line you DON\textquotesingle{}T want to commit}
\ExtensionTok{Stage}\NormalTok{ this hunk }\PreprocessorTok{[}\SpecialStringTok{y,n,q,a,d,s,e,?}\PreprocessorTok{]?}
\end{Highlighting}
\end{Shaded}

You answer \texttt{y} to the fix hunk and \texttt{n} to the \texttt{breakpoint()} hunk \(\rightarrow\) your commit contains only the fix. This is a \emph{hygiene} superpower: it lets you craft clean, single-purpose commits even when your working directory is messy. \texttt{s} splits a hunk; \texttt{e} lets you hand-edit which lines stage.

\begin{quote}
\textbf{Interview takeaway:} The index exists so a commit is an \emph{intentional, curated snapshot}, not "whatever happens to be on disk." Mentioning \texttt{git\ add\ -p} to keep commits atomic signals strong hygiene.
\end{quote}

\section{Branching \& Merging}\label{branching--merging}

\subsection{A branch is a ref; here\textquotesingle s what creating/switching does}\label{a-branch-is-a-ref-heres-what-creatingswitching-does}

\setcodelang{Bash}

\begin{Shaded}
\begin{Highlighting}[]
\FunctionTok{git}\NormalTok{ branch feature/login        }\CommentTok{\# write .git/refs/heads/feature/login = \textless{}current HEAD hash\textgreater{}}
\FunctionTok{git}\NormalTok{ switch feature/login        }\CommentTok{\# point HEAD at that branch (update working dir to match)}
\FunctionTok{git}\NormalTok{ switch }\AttributeTok{{-}c}\NormalTok{ feature/login     }\CommentTok{\# create AND switch in one step}
\end{Highlighting}
\end{Shaded}

Creating a branch is O(1): write one 41-byte file. Switching updates HEAD and reconciles your working tree to the target commit\textquotesingle s snapshot.

\subsection{Fast-forward merge (no divergence)}\label{fast-forward-merge-no-divergence}

If the branch you\textquotesingle re merging \emph{in} is strictly ahead of your current branch --- your branch is an ancestor of it --- Git can just \textbf{slide the pointer forward}. No merge commit, no new content.

\setcodelang{}

\begin{verbatim}
Before:                          After  git merge feature  (fast-forward):
  main ─▶ A ── B                   main and feature both ─▶ A ── B ── C ── D
                 \                                                         ▲
  feature ─▶      C ── D                                          (just moved the
                                                                   main pointer)
\end{verbatim}

\texttt{main} had no commits of its own that \texttt{feature} lacked, so merging is just \texttt{main\ =\ feature}. Use \texttt{git\ merge\ -\/-ff-only} to \emph{require} a fast-forward (fail loudly if a real merge would be needed) --- common in CI to keep history linear. Use \texttt{-\/-no-ff} to \emph{force} a merge commit even when a fast-forward is possible (preserves the "this was a feature branch" topology).

\subsection{Three-way merge (divergent branches)}\label{three-way-merge-divergent-branches}

When both branches have new commits since they split, Git performs a \textbf{3-way merge} using the \textbf{common ancestor} (merge base) plus the two tips:

\setcodelang{}

\begin{verbatim}
            ┌── C ── D ──┐  (feature: your work)
   A ── B ──┤             ├──▶ M  (merge commit, TWO parents: D and F)
            └── E ── F ──┘  (main: someone else's work)
            ▲
        merge base (B) = nearest common ancestor
\end{verbatim}

The algorithm (default strategy \textbf{ort}, formerly \textbf{recursive}):

\begin{enumerate}
\def\labelenumi{\arabic{enumi}.}
\tightlist
\item
  Find the \textbf{merge base} B (\texttt{git\ merge-base\ main\ feature}).
\item
  For each file, compare three versions: \textbf{base (B)}, \textbf{ours (main)}, \textbf{theirs (feature)}.
\item
  If only one side changed a region \(\rightarrow\) take that change automatically.
\item
  If both sides changed the \emph{same} region differently \(\rightarrow\) \textbf{conflict}; ask the human.
\item
  Produce a new \textbf{merge commit M} with \emph{two parents} (D and F), recording that both lines of history merged here.
\end{enumerate}

Why "3-way" beats "2-way": a plain 2-way diff between two files can\textquotesingle t tell \emph{which} side made a change --- it only sees they differ. The base gives Git the reference point to say "ours added this line; theirs is unchanged here, so keep ours" vs "both edited this same line \(\rightarrow\) conflict."

The \textbf{recursive/ort} strategy handles the case where two branches have \emph{multiple} common ancestors (criss-cross merges): it recursively merges the ancestors into a single virtual base. \textbf{ort} ("Ostensibly Recursive\textquotesingle s Twin") is a faster, more correct rewrite that became the default in Git 2.34 --- better rename detection, fewer spurious conflicts, much faster on big repos.

\subsection{Merge conflicts: markers and resolution}\label{merge-conflicts-markers-and-resolution}

When both sides edit the same lines, Git writes both versions into the file with markers and pauses:

\setcodelang{Python}

\begin{Shaded}
\begin{Highlighting}[]
\OperatorTok{\textless{}\textless{}\textless{}\textless{}\textless{}\textless{}\textless{}}\NormalTok{ HEAD (ours — current branch)}
\NormalTok{timeout }\OperatorTok{=} \DecValTok{30}
\OperatorTok{=======}
\NormalTok{timeout }\OperatorTok{=} \DecValTok{60}
\OperatorTok{\textgreater{}\textgreater{}\textgreater{}\textgreater{}\textgreater{}\textgreater{}\textgreater{}}\NormalTok{ feature}\OperatorTok{/}\NormalTok{tune}\OperatorTok{{-}}\NormalTok{timeouts (theirs — incoming)}
\end{Highlighting}
\end{Shaded}

\setcodelang{Bash}

\begin{Shaded}
\begin{Highlighting}[]
\FunctionTok{git}\NormalTok{ status                 }\CommentTok{\# lists "both modified" files}
\CommentTok{\# edit the file: pick one, combine them, delete ALL the \textless{}\textless{}\textless{}\textless{} ==== \textgreater{}\textgreater{}\textgreater{}\textgreater{} markers}
\FunctionTok{git}\NormalTok{ add config.py          }\CommentTok{\# mark THIS file resolved}
\FunctionTok{git}\NormalTok{ merge }\AttributeTok{{-}{-}continue}       \CommentTok{\# (or git commit) once all conflicts resolved}
\FunctionTok{git}\NormalTok{ merge }\AttributeTok{{-}{-}abort}          \CommentTok{\# bail out entirely, restore pre{-}merge state}

\CommentTok{\# Shortcuts when you know which whole side you want, per file:}
\FunctionTok{git}\NormalTok{ checkout }\AttributeTok{{-}{-}ours}\NormalTok{   config.py }\KeywordTok{\&\&} \FunctionTok{git}\NormalTok{ add config.py   }\CommentTok{\# keep current branch\textquotesingle{}s version}
\FunctionTok{git}\NormalTok{ checkout }\AttributeTok{{-}{-}theirs}\NormalTok{ config.py }\KeywordTok{\&\&} \FunctionTok{git}\NormalTok{ add config.py   }\CommentTok{\# keep incoming version}
\end{Highlighting}
\end{Shaded}

Enable \textbf{\texttt{diff3}/\texttt{zdiff3} conflict style} to \emph{also} show the merge base --- it makes resolution far easier because you can see what the original was:

\setcodelang{Bash}

\begin{Shaded}
\begin{Highlighting}[]
\ExtensionTok{$}\NormalTok{ git config }\AttributeTok{{-}{-}global}\NormalTok{ merge.conflictstyle zdiff3}
\OperatorTok{\textless{}\textless{}\textless{}\textless{}\textless{}\textless{}\textless{}}\NormalTok{ HEAD}
\FunctionTok{timeout}\NormalTok{ = 30}
\KeywordTok{|||||||} \ExtensionTok{base}
\FunctionTok{timeout}\NormalTok{ = 45         }\CommentTok{\# ← the original; now you can see WHO changed WHAT}
\ExtensionTok{=======}
\FunctionTok{timeout}\NormalTok{ = 60}
\OperatorTok{\textgreater{}\textgreater{}\textgreater{}\textgreater{}\textgreater{}\textgreater{}\textgreater{}}\NormalTok{ feature}
\end{Highlighting}
\end{Shaded}

Turn on \textbf{\texttt{rerere}} ("reuse recorded resolution") so Git remembers how you resolved a given conflict and auto-applies it next time the same conflict reappears (huge for long-lived branches and repeated rebases):

\setcodelang{Bash}

\begin{Shaded}
\begin{Highlighting}[]
\FunctionTok{git}\NormalTok{ config }\AttributeTok{{-}{-}global}\NormalTok{ rerere.enabled true}
\end{Highlighting}
\end{Shaded}

\subsection{Octopus merge (3+ branches at once)}\label{octopus-merge-3-branches-at-once}

\texttt{git\ merge\ a\ b\ c} creates a single commit with \emph{more than two} parents --- an \textbf{octopus merge}. It only works when there are no conflicts (it refuses if manual resolution is needed). Mostly used to bundle independent topic branches; rarely seen day-to-day, but a good "do you know edge cases?" answer.

\setcodelang{}

\begin{verbatim}
        ┌─ topicA ─┐
   main ┼─ topicB ─┼─▶ M (3 parents)
        └─ topicC ─┘
\end{verbatim}

\begin{quote}
\textbf{Interview takeaway:} Be able to draw the merge-base diagram and explain \emph{why} the third reference point (the base) is what makes auto-merging possible. Then explain that a conflict is simply "both sides changed the same region relative to the base, so Git refuses to guess."
\end{quote}

\section{Rebase vs Merge}\label{rebase-vs-merge}

Both integrate changes from one branch into another. The difference is \emph{topology} and \emph{whether new commits are created}.

\subsection{Merge = preserve history, add a merge commit}\label{merge--preserve-history-add-a-merge-commit}

Merge keeps the true, branching history and ties it together with a merge commit (two parents). Nothing is rewritten; all original commits keep their hashes.

\subsection{Rebase = replay your commits onto a new base (new hashes!)}\label{rebase--replay-your-commits-onto-a-new-base-new-hashes}

\texttt{git\ rebase\ main} takes \emph{your} commits, sets them aside, fast-forwards your branch to \texttt{main}, then \textbf{re-applies each of your commits one at a time} on top --- producing \textbf{brand-new commits with new hashes} (different parents \(\rightarrow\) different commit text \(\rightarrow\) different SHA). The result is a clean linear history \emph{as if} you\textquotesingle d started your work from the latest main.

\setcodelang{}

\begin{verbatim}
Before rebase:                          After  git rebase main  (on feature):
            C ── D  (feature)                            C'── D'  (feature, NEW hashes)
           /                                            /
  A ── B ──                                A ── B ── E ── F
           \                                   (main)
            E ── F  (main)

  C,D had parent B.                      C',D' now have parent F. Same DIFFS,
                                          new COMMITS. C/D become dangling.
\end{verbatim}

\setcodelang{Bash}

\begin{Shaded}
\begin{Highlighting}[]
\FunctionTok{git}\NormalTok{ switch feature}
\FunctionTok{git}\NormalTok{ rebase main          }\CommentTok{\# replay feature\textquotesingle{}s commits on top of main}
\CommentTok{\# conflicts? resolve, then:}
\FunctionTok{git}\NormalTok{ add }\OperatorTok{\textless{}}\NormalTok{file}\OperatorTok{\textgreater{}} \KeywordTok{\&\&} \FunctionTok{git}\NormalTok{ rebase }\AttributeTok{{-}{-}continue}
\FunctionTok{git}\NormalTok{ rebase }\AttributeTok{{-}{-}abort}       \CommentTok{\# back out, restore original state}
\FunctionTok{git}\NormalTok{ rebase }\AttributeTok{{-}{-}skip}        \CommentTok{\# drop the current conflicting commit and continue}
\end{Highlighting}
\end{Shaded}

\subsection{Side-by-side}\label{side-by-side}

\begin{longtable}[]{@{}
  >{\raggedright\arraybackslash}p{(\columnwidth - 4\tabcolsep) * \real{0.3608}}
  >{\raggedright\arraybackslash}p{(\columnwidth - 4\tabcolsep) * \real{0.2990}}
  >{\raggedright\arraybackslash}p{(\columnwidth - 4\tabcolsep) * \real{0.3402}}@{}}
\toprule\noalign{}
\rowcolor{acc!16}
\begin{minipage}[b]{\linewidth}\raggedright
Aspect
\end{minipage} & \begin{minipage}[b]{\linewidth}\raggedright
Merge
\end{minipage} & \begin{minipage}[b]{\linewidth}\raggedright
Rebase
\end{minipage} \\
\midrule\noalign{}
\endhead
\bottomrule\noalign{}
\endlastfoot
History shape & Branching, with merge commits & Linear, flat \\
\rowcolor{zebra}
Rewrites commits? & No (original hashes kept) & \textbf{Yes} (new hashes) \\
Traceability of "when branches met" & Preserved (merge commit) & Lost (looks sequential) \\
\rowcolor{zebra}
Conflict resolution & Once, at the merge & Possibly once \emph{per replayed commit} \\
Safe on shared/pushed branches? & Yes & \textbf{No} (see Golden Rule) \\
\rowcolor{zebra}
\texttt{git\ bisect} / \texttt{git\ log} readability & Noisier graph & Cleaner, easier to bisect \\
\end{longtable}

\subsection{Interactive rebase --- the history-editing superpower}\label{interactive-rebase--the-history-editing-superpower}

\texttt{git\ rebase\ -i} opens an editor listing your commits with an action per line. This is how you craft a clean PR out of messy work-in-progress commits.

\setcodelang{Bash}

\begin{Shaded}
\begin{Highlighting}[]
\ExtensionTok{$}\NormalTok{ git rebase }\AttributeTok{{-}i}\NormalTok{ HEAD\textasciitilde{}5}
\ExtensionTok{pick}\NormalTok{   a1b2c3 Add login form}
\ExtensionTok{squash}\NormalTok{ 4d5e6f WIP fix typo            }\CommentTok{\# fold into previous, combine messages}
\ExtensionTok{fixup}\NormalTok{  7g8h9i oops forgot import      }\CommentTok{\# fold into previous, DISCARD this message}
\ExtensionTok{reword}\NormalTok{ 0j1k2l Add validatoin          }\CommentTok{\# keep commit, edit its message (fix typo)}
\ExtensionTok{edit}\NormalTok{   3m4n5o Refactor auth service   }\CommentTok{\# stop here so you can amend/split}
\ExtensionTok{drop}\NormalTok{   6p7q8r debug print statement   }\CommentTok{\# delete this commit entirely}
\CommentTok{\# you can also REORDER lines to reorder commits}
\end{Highlighting}
\end{Shaded}

\begin{longtable}[]{@{}
  >{\raggedright\arraybackslash}p{(\columnwidth - 2\tabcolsep) * \real{0.2394}}
  >{\raggedright\arraybackslash}p{(\columnwidth - 2\tabcolsep) * \real{0.7606}}@{}}
\toprule\noalign{}
\rowcolor{acc!16}
\begin{minipage}[b]{\linewidth}\raggedright
Action
\end{minipage} & \begin{minipage}[b]{\linewidth}\raggedright
Effect
\end{minipage} \\
\midrule\noalign{}
\endhead
\bottomrule\noalign{}
\endlastfoot
\texttt{pick} & keep the commit as-is \\
\rowcolor{zebra}
\texttt{reword} & keep the commit, edit its message \\
\texttt{edit} & pause at this commit to amend it or split it into several \\
\rowcolor{zebra}
\texttt{squash} (\texttt{s}) & merge into the previous commit, \emph{combine} both messages \\
\texttt{fixup} (\texttt{f}) & merge into the previous commit, \emph{discard} this message \\
\rowcolor{zebra}
\texttt{drop} (\texttt{d}) & remove the commit entirely \\
(reorder lines) & reorders the commits \\
\end{longtable}

\textbf{Worked example --- turning 5 messy commits into 2 clean ones:}

\setcodelang{}

\begin{verbatim}
WIP commits in your branch:               After interactive rebase:
  a1b2  Add login form                      X1  Add login form + validation
  4d5e  WIP fix typo            ─squash─▶        (3 commits folded into one,
  7g8h  oops forgot import      ─fixup──▶         message cleaned)
  0j1k  Add validatoin          ─reword─▶
  6p7q  debug print statement   ─drop───▶    (gone)
\end{verbatim}

You end up with two logically distinct, well-described commits a reviewer will thank you for.

\subsection{Autosquash --- fixups without manual reordering}\label{autosquash--fixups-without-manual-reordering}

Make a fix and tag it to a specific earlier commit; later, autosquash sorts and folds it automatically:

\setcodelang{Bash}

\begin{Shaded}
\begin{Highlighting}[]
\FunctionTok{git}\NormalTok{ commit }\AttributeTok{{-}{-}fixup}\OperatorTok{=}\NormalTok{a1b2c3        }\CommentTok{\# creates "fixup! Add login form"}
\CommentTok{\# ...more work...}
\FunctionTok{git}\NormalTok{ rebase }\AttributeTok{{-}i} \AttributeTok{{-}{-}autosquash}\NormalTok{ main  }\CommentTok{\# Git auto{-}positions the fixup under a1b2c3 and marks it \textquotesingle{}fixup\textquotesingle{}}
\end{Highlighting}
\end{Shaded}

Set \texttt{git\ config\ -\/-global\ rebase.autosquash\ true} to make \texttt{-\/-autosquash} the default.

\subsection{The GOLDEN RULE: never rebase commits that have been pushed/shared}\label{the-golden-rule-never-rebase-commits-that-have-been-pushedshared}

\textbf{Rule:} \emph{Only rebase commits that live solely in your local repo and that nobody else has based work on.} Never rebase a branch others have pulled (like \texttt{main}, \texttt{develop}, or a shared feature branch).

\textbf{Exactly why it breaks collaborators:} rebasing replaces commits \texttt{C,\ D} with new commits \texttt{C\textquotesingle{},\ D\textquotesingle{}} (different hashes). Suppose you already pushed \texttt{C,\ D} and a teammate pulled them and built \texttt{G} on top of \texttt{D}. Now you rebase, force-push \texttt{C\textquotesingle{},\ D\textquotesingle{}}, and:

\setcodelang{}

\begin{verbatim}
Remote now has:  A ── B ── E ── F ── C'── D'
Teammate has:    A ── B ──────────── C ── D ── G   (built on the OLD D)
\end{verbatim}

When your teammate pulls, Git sees their \texttt{C,\ D} as \emph{unrelated} to your \texttt{C\textquotesingle{},\ D\textquotesingle{}} (different hashes). Their history and yours have \textbf{diverged}. They get a messy merge, duplicate commits, and conflicts --- and if they force-push to "fix" it, they may clobber \emph{your} \texttt{C\textquotesingle{},\ D\textquotesingle{}}. The old commits they depended on effectively vanished from the shared branch. Multiply across a team and you get "the repo is broken" Slack threads. The fix-up dance is painful (\texttt{git\ rebase\ -\/-onto}), which is why the rule is \emph{don\textquotesingle t}.

\subsection{\texorpdfstring{\texttt{git\ pull\ -\/-rebase}}{git pull -\/-rebase}}\label{git-pull---rebase}

Default \texttt{git\ pull} = \texttt{fetch} + \textbf{merge}, which sprinkles tiny "Merge branch \textquotesingle main\textquotesingle{} of origin" commits all over history. \texttt{git\ pull\ -\/-rebase} = \texttt{fetch} + \textbf{rebase} your local commits on top of the fetched remote tip \(\rightarrow\) linear history, no noise. Safe here because your \emph{local-only} commits are the ones being replayed (not shared ones). Many teams set it as default:

\setcodelang{Bash}

\begin{Shaded}
\begin{Highlighting}[]
\FunctionTok{git}\NormalTok{ config }\AttributeTok{{-}{-}global}\NormalTok{ pull.rebase true}
\CommentTok{\# or per{-}pull:}
\FunctionTok{git}\NormalTok{ pull }\AttributeTok{{-}{-}rebase}\NormalTok{ origin main}
\end{Highlighting}
\end{Shaded}

\begin{quote}
\textbf{Interview takeaway:} "Rebase rewrites commits into new hashes for a linear history; merge preserves topology with a merge commit. Rebase locally for cleanliness, merge to integrate shared branches, and \emph{never} rebase anything already pushed and depended upon." That one sentence is the senior-level answer.
\end{quote}

\section{Undo, Reset \& Recovery}\label{undo-reset--recovery}

The thing that trips everyone up: \texttt{reset}, \texttt{revert}, \texttt{restore}, and \texttt{checkout} all "undo," but touch \emph{different areas} with \emph{different safety profiles}.

\subsection{\texorpdfstring{\texttt{git\ reset\ -\/-soft\ /\ -\/-mixed\ /\ -\/-hard} --- what each touches}{git reset -\/-soft / -\/-mixed / -\/-hard --- what each touches}}\label{git-reset---soft----mixed----hard--what-each-touches}

\texttt{reset} moves the current \textbf{branch pointer} to a target commit and \emph{optionally} updates the index and working dir. The mode controls how far the changes "fall back":

\begin{longtable}[]{@{}
  >{\raggedright\arraybackslash}p{(\columnwidth - 8\tabcolsep) * \real{0.1244}}
  >{\raggedright\arraybackslash}p{(\columnwidth - 8\tabcolsep) * \real{0.1430}}
  >{\raggedright\arraybackslash}p{(\columnwidth - 8\tabcolsep) * \real{0.1010}}
  >{\raggedright\arraybackslash}p{(\columnwidth - 8\tabcolsep) * \real{0.1514}}
  >{\raggedright\arraybackslash}p{(\columnwidth - 8\tabcolsep) * \real{0.4802}}@{}}
\toprule\noalign{}
\rowcolor{acc!16}
\begin{minipage}[b]{\linewidth}\raggedright
Mode
\end{minipage} & \begin{minipage}[b]{\linewidth}\raggedright
Moves branch HEAD
\end{minipage} & \begin{minipage}[b]{\linewidth}\raggedright
Resets \textbf{Index}
\end{minipage} & \begin{minipage}[b]{\linewidth}\raggedright
Resets \textbf{Working Dir}
\end{minipage} & \begin{minipage}[b]{\linewidth}\raggedright
Result / use
\end{minipage} \\
\midrule\noalign{}
\endhead
\bottomrule\noalign{}
\endlastfoot
\texttt{-\/-soft} & & & & Changes from undone commits become \textbf{staged}. Use to re-commit differently / combine last N commits. \\
\rowcolor{zebra}
\texttt{-\/-mixed} (default) & & & & Changes become \textbf{unstaged} edits in working dir. Use to unstage / redo the commit. \\
\texttt{-\/-hard} & & & & \textbf{Everything wiped} to target. \emph{Destructive} --- uncommitted work is gone. \\
\end{longtable}

\setcodelang{}

\begin{verbatim}
                 ┌─ working dir
        ┌─ index ┤
   HEAD ┤        └─ working dir
   --soft : move HEAD only          → diffs sit in INDEX (staged)
   --mixed: move HEAD + index       → diffs sit in WORKING DIR (unstaged)
   --hard : move HEAD + index + WD  → diffs DELETED
\end{verbatim}

\setcodelang{Bash}

\begin{Shaded}
\begin{Highlighting}[]
\FunctionTok{git}\NormalTok{ reset }\AttributeTok{{-}{-}soft}\NormalTok{ HEAD\textasciitilde{}3   }\CommentTok{\# "undo last 3 commits but keep all changes staged" (squash{-}by{-}hand)}
\FunctionTok{git}\NormalTok{ reset }\AttributeTok{{-}{-}mixed}\NormalTok{ HEAD\textasciitilde{}1  }\CommentTok{\# "undo last commit, keep changes as edits to re{-}stage"}
\FunctionTok{git}\NormalTok{ reset }\AttributeTok{{-}{-}hard}\NormalTok{ HEAD\textasciitilde{}1   }\CommentTok{\# "throw away last commit AND its changes" — danger}
\FunctionTok{git}\NormalTok{ reset }\AttributeTok{{-}{-}hard}\NormalTok{ origin/main  }\CommentTok{\# "make my branch exactly match the remote" — danger}
\end{Highlighting}
\end{Shaded}

\begin{quote}
\texttt{-\/-hard} is the only common Git command that destroys \emph{uncommitted} working-tree changes irrecoverably (they were never in a commit, so reflog can\textquotesingle t help). Committed work it removes is still recoverable via reflog (below).
\end{quote}

\subsection{\texorpdfstring{\texttt{git\ revert} --- the safe, public undo}{git revert --- the safe, public undo}}\label{git-revert--the-safe-public-undo}

\texttt{revert} does \textbf{not} rewrite history. It creates a \textbf{new commit} that applies the \emph{inverse} of a target commit\textquotesingle s diff. History moves \emph{forward}; the bad change is neutralized without removing anything. This is the \textbf{only} correct way to undo a commit on a \textbf{shared/public} branch.

\setcodelang{}

\begin{verbatim}
Before:  A ── B ── C(bad) ── D
After  git revert C:
         A ── B ── C(bad) ── D ── C⁻¹   (new commit that undoes C; C still in history)
\end{verbatim}

\setcodelang{Bash}

\begin{Shaded}
\begin{Highlighting}[]
\FunctionTok{git}\NormalTok{ revert }\OperatorTok{\textless{}}\NormalTok{bad{-}sha}\OperatorTok{\textgreater{}}          \CommentTok{\# make an inverse commit (opens editor for message)}
\FunctionTok{git}\NormalTok{ revert HEAD               }\CommentTok{\# undo the most recent commit, safely}
\FunctionTok{git}\NormalTok{ revert }\AttributeTok{{-}m}\NormalTok{ 1 }\OperatorTok{\textless{}}\NormalTok{merge{-}sha}\OperatorTok{\textgreater{}}   \CommentTok{\# revert a MERGE commit ({-}m 1 = keep first parent\textquotesingle{}s line)}
\FunctionTok{git}\NormalTok{ revert }\AttributeTok{{-}{-}no{-}commit}\NormalTok{ A B C  }\CommentTok{\# stage inverses of several commits into one revert}
\end{Highlighting}
\end{Shaded}

\subsection{\texorpdfstring{\texttt{reset} vs \texttt{revert} --- when each is safe}{reset vs revert --- when each is safe}}\label{reset-vs-revert--when-each-is-safe}

\begin{longtable}[]{@{}
  >{\raggedright\arraybackslash}p{(\columnwidth - 4\tabcolsep) * \real{0.2200}}
  >{\raggedright\arraybackslash}p{(\columnwidth - 4\tabcolsep) * \real{0.4000}}
  >{\raggedright\arraybackslash}p{(\columnwidth - 4\tabcolsep) * \real{0.3800}}@{}}
\toprule\noalign{}
\rowcolor{acc!16}
\begin{minipage}[b]{\linewidth}\raggedright
\end{minipage} & \begin{minipage}[b]{\linewidth}\raggedright
\texttt{git\ reset}
\end{minipage} & \begin{minipage}[b]{\linewidth}\raggedright
\texttt{git\ revert}
\end{minipage} \\
\midrule\noalign{}
\endhead
\bottomrule\noalign{}
\endlastfoot
Mechanism & Moves branch pointer back (rewrites tip) & Adds an inverse commit (moves forward) \\
\rowcolor{zebra}
Rewrites history? & Yes & No \\
Safe on shared branch? & \textbf{No} (forces force-push) & \textbf{Yes} \\
\rowcolor{zebra}
Use when & Local, unpushed cleanup & Undoing something already pushed \\
\end{longtable}

\subsection{\texorpdfstring{\texttt{restore} / \texttt{checkout} --- file-level}{restore / checkout --- file-level}}\label{restore--checkout--file-level}

\setcodelang{Bash}

\begin{Shaded}
\begin{Highlighting}[]
\FunctionTok{git}\NormalTok{ restore file.py                 }\CommentTok{\# discard unstaged edits to file (from index/HEAD) — destructive}
\FunctionTok{git}\NormalTok{ restore }\AttributeTok{{-}{-}staged}\NormalTok{ file.py        }\CommentTok{\# unstage (index → keep working edits)}
\FunctionTok{git}\NormalTok{ restore }\AttributeTok{{-}{-}source}\OperatorTok{=}\NormalTok{HEAD\textasciitilde{}2 app.py  }\CommentTok{\# bring an OLD version of a file into working dir}
\FunctionTok{git}\NormalTok{ checkout }\OperatorTok{\textless{}}\NormalTok{commit}\OperatorTok{\textgreater{}}\NormalTok{ {-}{-} file.py    }\CommentTok{\# older equivalent of restore {-}{-}source}
\end{Highlighting}
\end{Shaded}

\subsection{\texorpdfstring{\texttt{git\ clean} --- remove untracked files}{git clean --- remove untracked files}}\label{git-clean--remove-untracked-files}

Git won\textquotesingle t auto-delete \emph{untracked} files; \texttt{clean} does, and it\textquotesingle s destructive (no reflog for files never tracked). \textbf{Always dry-run first.}

\setcodelang{Bash}

\begin{Shaded}
\begin{Highlighting}[]
\FunctionTok{git}\NormalTok{ clean }\AttributeTok{{-}n}         \CommentTok{\# DRY RUN: show what WOULD be deleted}
\FunctionTok{git}\NormalTok{ clean }\AttributeTok{{-}fd}        \CommentTok{\# actually delete untracked files ({-}f) and dirs ({-}d)}
\FunctionTok{git}\NormalTok{ clean }\AttributeTok{{-}fdx}       \CommentTok{\# also delete ignored files (build artifacts, node\_modules) — be careful}
\end{Highlighting}
\end{Shaded}

\subsection{Safety summary}\label{safety-summary}

\setcodelang{}

\begin{verbatim}
Recoverable via reflog (commits were made):  amend, reset (any), rebase, branch delete
NOT recoverable (never committed):           reset --hard on uncommitted edits,
                                             git clean, git restore <file> (unstaged edits),
                                             git stash drop (after the fact)
\end{verbatim}

\begin{quote}
\textbf{Interview takeaway:} "\texttt{reset} rewinds the branch pointer (and optionally index/working dir) --- local use only; \texttt{revert} adds an inverse commit and is the safe way to undo something already pushed. \texttt{-\/-hard} and \texttt{clean} are the two commands that can lose \emph{uncommitted} work for good."
\end{quote}

\section{Powerful Commands (bisect, cherry-pick, reflog, stash, worktree)}\label{powerful-commands-bisect-cherry-pick-reflog-stash-worktree}

\subsection{\texorpdfstring{\texttt{git\ reflog} --- your undo history / safety net}{git reflog --- your undo history / safety net}}\label{git-reflog--your-undo-history--safety-net}

The \textbf{reflog} records \emph{every} movement of HEAD and branch tips locally (commits, resets, rebases, checkouts, merges) --- even ones that left commits "dangling" and unreachable from any branch. It\textquotesingle s how you recover from almost any "I lost my work" disaster. Reflog entries are local, not pushed, and expire after \textasciitilde90 days (30 for unreachable).

\setcodelang{Bash}

\begin{Shaded}
\begin{Highlighting}[]
\ExtensionTok{$}\NormalTok{ git reflog}
\ExtensionTok{3f8a1c9}\NormalTok{ HEAD@\{0\}: reset: moving to HEAD\textasciitilde{}3}
\ExtensionTok{a1b2c3d}\NormalTok{ HEAD@\{1\}: commit: Add payment retry}
\ExtensionTok{7c4b2ae}\NormalTok{ HEAD@\{2\}: commit: Add login handler}
\ExtensionTok{e2f3a4b}\NormalTok{ HEAD@\{3\}: checkout: moving from main to feature}
\end{Highlighting}
\end{Shaded}

\textbf{Scenario A --- undo a bad \texttt{reset\ -\/-hard}:} You ran \texttt{git\ reset\ -\/-hard\ HEAD\textasciitilde{}3} and panic. The three commits are unreachable but not gone yet.

\setcodelang{Bash}

\begin{Shaded}
\begin{Highlighting}[]
\FunctionTok{git}\NormalTok{ reflog                       }\CommentTok{\# find the hash BEFORE the reset (e.g. a1b2c3d HEAD@\{1\})}
\FunctionTok{git}\NormalTok{ reset }\AttributeTok{{-}{-}hard}\NormalTok{ a1b2c3d         }\CommentTok{\# restore the branch to that exact state}
\end{Highlighting}
\end{Shaded}

\textbf{Scenario B --- recover a deleted branch:}

\setcodelang{Bash}

\begin{Shaded}
\begin{Highlighting}[]
\FunctionTok{git}\NormalTok{ branch }\AttributeTok{{-}D}\NormalTok{ feature/login      }\CommentTok{\# oops, deleted unmerged branch}
\FunctionTok{git}\NormalTok{ reflog                       }\CommentTok{\# find the last commit that WAS feature/login\textquotesingle{}s tip}
\FunctionTok{git}\NormalTok{ switch }\AttributeTok{{-}c}\NormalTok{ feature/login }\OperatorTok{\textless{}}\NormalTok{that{-}sha}\OperatorTok{\textgreater{}}   \CommentTok{\# recreate the branch at that commit}
\end{Highlighting}
\end{Shaded}

\textbf{Scenario C --- undo a botched rebase:} Before any rebase, \texttt{HEAD@\{n\}} records the pre-rebase tip.

\setcodelang{Bash}

\begin{Shaded}
\begin{Highlighting}[]
\FunctionTok{git}\NormalTok{ reflog                       }\CommentTok{\# find "rebase (start): ..." — the commit just before it is your old tip}
\FunctionTok{git}\NormalTok{ reset }\AttributeTok{{-}{-}hard}\NormalTok{ HEAD@\{5\}        }\CommentTok{\# or the specific pre{-}rebase hash}
\CommentTok{\# Even simpler if it was the last operation:}
\FunctionTok{git}\NormalTok{ reset }\AttributeTok{{-}{-}hard}\NormalTok{ ORIG\_HEAD       }\CommentTok{\# ORIG\_HEAD = where HEAD was before the last "big" move}
\end{Highlighting}
\end{Shaded}

\subsection{\texorpdfstring{\texttt{git\ bisect} --- binary-search a regression in O(log n)}{git bisect --- binary-search a regression in O(log n)}}\label{git-bisect--binary-search-a-regression-in-olog-n}

A test passes at an old commit and fails now. Somewhere in the (say) 1,024 commits between, one broke it. Linear search = up to 1,024 checks. Bisect = \texttt{log2(1024)} = \textbf{\textasciitilde10 checks} because each step halves the suspect range.

\setcodelang{Bash}

\begin{Shaded}
\begin{Highlighting}[]
\FunctionTok{git}\NormalTok{ bisect start}
\FunctionTok{git}\NormalTok{ bisect bad                 }\CommentTok{\# current commit is broken}
\FunctionTok{git}\NormalTok{ bisect good v1.4.0         }\CommentTok{\# this old tag was fine}
\CommentTok{\# Git checks out the MIDPOINT commit. You test it, then tell Git:}
\FunctionTok{git}\NormalTok{ bisect good                }\CommentTok{\#   ...if this midpoint works}
\FunctionTok{git}\NormalTok{ bisect bad                 }\CommentTok{\#   ...if it\textquotesingle{}s broken}
\CommentTok{\# repeat \textasciitilde{}log2(N) times; Git narrows to the FIRST bad commit}
\FunctionTok{git}\NormalTok{ bisect reset               }\CommentTok{\# done — return to your original HEAD}
\end{Highlighting}
\end{Shaded}

\textbf{Automate it} with a script that exits 0 (good) / non-zero (bad):

\setcodelang{Bash}

\begin{Shaded}
\begin{Highlighting}[]
\FunctionTok{git}\NormalTok{ bisect start HEAD v1.4.0   }\CommentTok{\# bad=HEAD, good=v1.4.0 in one line}
\FunctionTok{git}\NormalTok{ bisect run pytest tests/test\_checkout.py::test\_total   }\CommentTok{\# Git runs it at each step automatically}
\CommentTok{\# Use exit code 125 in your script to say "skip — can\textquotesingle{}t test this commit"}
\end{Highlighting}
\end{Shaded}

The output \texttt{\textless{}sha\textgreater{}\ is\ the\ first\ bad\ commit} plus that commit\textquotesingle s diff usually points straight at the bug. Bisect is the single biggest argument for keeping commits \textbf{small and each-one-building} (so you can bisect cleanly).

\subsection{\texorpdfstring{\texttt{git\ cherry-pick} --- copy a commit onto another branch}{git cherry-pick --- copy a commit onto another branch}}\label{git-cherry-pick--copy-a-commit-onto-another-branch}

Applies the \emph{diff} of a specific commit as a \textbf{new commit} (new hash) on your current branch. The canonical use: a bug is fixed on \texttt{main}, and you need that exact fix on the \texttt{release/2.3} branch without pulling in everything else.

\setcodelang{Bash}

\begin{Shaded}
\begin{Highlighting}[]
\FunctionTok{git}\NormalTok{ switch release/2.3}
\FunctionTok{git}\NormalTok{ cherry{-}pick a1b2c3d              }\CommentTok{\# apply that fix here as a new commit}
\FunctionTok{git}\NormalTok{ cherry{-}pick a1b2c3d\^{}..f6e7d8     }\CommentTok{\# a RANGE of commits}
\FunctionTok{git}\NormalTok{ cherry{-}pick }\AttributeTok{{-}x}\NormalTok{ a1b2c3d           }\CommentTok{\# append "(cherry picked from commit a1b2c3d)" to the message}
\CommentTok{\# conflict? resolve, then:}
\FunctionTok{git}\NormalTok{ add }\OperatorTok{\textless{}}\NormalTok{file}\OperatorTok{\textgreater{}} \KeywordTok{\&\&} \FunctionTok{git}\NormalTok{ cherry{-}pick }\AttributeTok{{-}{-}continue}
\FunctionTok{git}\NormalTok{ cherry{-}pick }\AttributeTok{{-}{-}abort}
\end{Highlighting}
\end{Shaded}

Caveat: cherry-picking the \emph{same} change onto multiple branches creates \emph{duplicate} commits (different hashes, same diff). When those branches later merge, Git is usually smart enough to detect the patch is already present, but it can cause spurious conflicts. Prefer cherry-pick for genuine hotfix backports, not as a substitute for merging.

\subsection{\texorpdfstring{\texttt{git\ stash} --- shelve work-in-progress}{git stash --- shelve work-in-progress}}\label{git-stash--shelve-work-in-progress}

Saves your uncommitted changes (working dir + staged) onto a stack and reverts to a clean tree, so you can switch context (e.g., urgent hotfix) without committing half-done work.

\setcodelang{Bash}

\begin{Shaded}
\begin{Highlighting}[]
\FunctionTok{git}\NormalTok{ stash                       }\CommentTok{\# shelve tracked changes, clean the working dir}
\FunctionTok{git}\NormalTok{ stash push }\AttributeTok{{-}m} \StringTok{"wip: auth"}   \CommentTok{\# named stash}
\FunctionTok{git}\NormalTok{ stash }\AttributeTok{{-}u}                    \CommentTok{\# also stash UNtracked files}
\FunctionTok{git}\NormalTok{ stash list                  }\CommentTok{\# stash@\{0\}: WIP on feature: ...}
\FunctionTok{git}\NormalTok{ stash show }\AttributeTok{{-}p}\NormalTok{ stash@\{0\}     }\CommentTok{\# view the diff}
\FunctionTok{git}\NormalTok{ stash pop                   }\CommentTok{\# re{-}apply the latest stash AND remove it from the stack}
\FunctionTok{git}\NormalTok{ stash apply stash@\{1\}       }\CommentTok{\# re{-}apply a specific stash, KEEP it on the stack}
\FunctionTok{git}\NormalTok{ stash branch fix stash@\{0\}  }\CommentTok{\# create a branch from a stash (if it won\textquotesingle{}t apply cleanly)}
\FunctionTok{git}\NormalTok{ stash drop stash@\{0\}        }\CommentTok{\# delete one stash}
\FunctionTok{git}\NormalTok{ stash clear                 }\CommentTok{\# nuke all stashes (no easy recovery)}
\end{Highlighting}
\end{Shaded}

Stashes are real commits under the hood (\texttt{refs/stash}), so a \emph{dropped} stash is recoverable via \texttt{git\ fsck\ -\/-unreachable\ \textbar{}\ grep\ commit} for a while --- but it\textquotesingle s fragile; don\textquotesingle t rely on it.

\subsection{\texorpdfstring{\texttt{git\ worktree} --- multiple working dirs, one repo}{git worktree --- multiple working dirs, one repo}}\label{git-worktree--multiple-working-dirs-one-repo}

Normally one repo = one working directory on one branch. \texttt{worktree} lets you check out \textbf{multiple branches simultaneously} in separate folders, all sharing the same \texttt{.git} object store (no second clone, no duplicated history). Great for: building a release branch while developing on a feature, or reviewing a PR without stashing.

\setcodelang{Bash}

\begin{Shaded}
\begin{Highlighting}[]
\FunctionTok{git}\NormalTok{ worktree add ../hotfix release/2.3   }\CommentTok{\# new folder ../hotfix checked out to release/2.3}
\FunctionTok{git}\NormalTok{ worktree add }\AttributeTok{{-}b}\NormalTok{ experiment ../exp    }\CommentTok{\# create branch \textquotesingle{}experiment\textquotesingle{} in ../exp}
\FunctionTok{git}\NormalTok{ worktree list}
\FunctionTok{git}\NormalTok{ worktree remove ../hotfix}
\end{Highlighting}
\end{Shaded}

\subsection{\texorpdfstring{\texttt{git\ blame} --- who/when/why for each line}{git blame --- who/when/why for each line}}\label{git-blame--whowhenwhy-for-each-line}

\setcodelang{Bash}

\begin{Shaded}
\begin{Highlighting}[]
\FunctionTok{git}\NormalTok{ blame app.py                  }\CommentTok{\# annotate every line with commit/author/date}
\FunctionTok{git}\NormalTok{ blame }\AttributeTok{{-}L}\NormalTok{ 40,60 app.py         }\CommentTok{\# only lines 40–60}
\FunctionTok{git}\NormalTok{ blame }\AttributeTok{{-}w} \AttributeTok{{-}C}\NormalTok{ app.py            }\CommentTok{\# ignore whitespace ({-}w), detect moved/copied code ({-}C)}
\FunctionTok{git}\NormalTok{ log }\AttributeTok{{-}S} \StringTok{"functionName"}         \CommentTok{\# "pickaxe": find commits that ADDED/REMOVED that string}
\FunctionTok{git}\NormalTok{ log }\AttributeTok{{-}L}\NormalTok{ :handler:app.py        }\CommentTok{\# follow the history of a single function}
\end{Highlighting}
\end{Shaded}

\texttt{blame} tells you the commit; the commit message tells you \emph{why} --- which is the whole point of writing good commit messages.

\subsection{\texorpdfstring{\texttt{git\ log} tricks}{git log tricks}}\label{git-log-tricks}

\setcodelang{Bash}

\begin{Shaded}
\begin{Highlighting}[]
\FunctionTok{git}\NormalTok{ log }\AttributeTok{{-}{-}oneline} \AttributeTok{{-}{-}graph} \AttributeTok{{-}{-}all} \AttributeTok{{-}{-}decorate}   \CommentTok{\# the classic visual history graph}
\FunctionTok{git}\NormalTok{ log }\AttributeTok{{-}{-}since}\OperatorTok{=}\StringTok{"2 weeks ago"} \AttributeTok{{-}{-}author}\OperatorTok{=}\StringTok{"Asha"}
\FunctionTok{git}\NormalTok{ log main..feature            }\CommentTok{\# commits on feature NOT yet on main (what your PR adds)}
\FunctionTok{git}\NormalTok{ log feature..main            }\CommentTok{\# commits on main that feature is missing}
\FunctionTok{git}\NormalTok{ log }\AttributeTok{{-}p}\NormalTok{ app.py                }\CommentTok{\# full diffs touching app.py}
\FunctionTok{git}\NormalTok{ shortlog }\AttributeTok{{-}sn}                 \CommentTok{\# commit counts per author}
\FunctionTok{git}\NormalTok{ log }\AttributeTok{{-}{-}first{-}parent}           \CommentTok{\# follow only mainline (skip merged{-}in branch detail)}
\end{Highlighting}
\end{Shaded}

\begin{quote}
\textbf{Interview takeaway:} Name-drop the \emph{right} tool for the \emph{right} problem: "regression \(\rightarrow\) \texttt{bisect}; lost commit \(\rightarrow\) \texttt{reflog}; backport a fix \(\rightarrow\) \texttt{cherry-pick}; context-switch \(\rightarrow\) \texttt{stash}; parallel branches \(\rightarrow\) \texttt{worktree}." That fluency reads as senior.
\end{quote}

\section{Branching Workflows}\label{branching-workflows}

A \emph{workflow} is the team agreement about how branches map to releases. The right choice depends on release cadence, team size, and CI maturity.

\subsection{Gitflow (heavyweight, release-train teams)}\label{gitflow-heavyweight-release-train-teams}

Two long-lived branches (\texttt{main} = production, \texttt{develop} = integration) plus short-lived \texttt{feature/*}, \texttt{release/*}, and \texttt{hotfix/*} branches.

\setcodelang{}

\begin{verbatim}
hotfix/* ───────────────────────────────●──▶ (merge to main AND develop, tag)
                                        /
main      ●────────────────●──────────●────────●  (every commit is a release, tagged)
           \              / \        /          \
release/*   \            ●───●  (stabilize: only bugfixes, version bump)
             \          /
develop   ●───●───●───●───●───●───●───●  (integration branch)
           \   \     /   /
feature/*   ●───●   ●───●   (branch off develop, merge back to develop)
\end{verbatim}

\begin{itemize}
\tightlist
\item
  \textbf{feature/} \(\rightarrow\) branch from \texttt{develop}, merge back to \texttt{develop}.
\item
  \textbf{release/} \(\rightarrow\) branch from \texttt{develop} to stabilize a version; only bugfixes; merge to \emph{both} \texttt{main} and \texttt{develop}; tag on \texttt{main}.
\item
  \textbf{hotfix/} \(\rightarrow\) branch from \texttt{main} for emergency prod fixes; merge to \emph{both} \texttt{main} and \texttt{develop}.
\end{itemize}

Good for: versioned software with scheduled releases (installed apps, libraries, mobile with app-store review). Downsides: many long-lived branches, painful merges, slow integration ("merge hell") --- overkill for continuously-deployed web services.

\subsection{GitHub Flow (lightweight, continuous deploy)}\label{github-flow-lightweight-continuous-deploy}

One long-lived branch (\texttt{main}, always deployable). Everything else is a short-lived feature branch \(\rightarrow\) PR \(\rightarrow\) review \(\rightarrow\) CI \(\rightarrow\) merge \(\rightarrow\) deploy.

\setcodelang{}

\begin{verbatim}
main  ●───────●───────●───────●  (always deployable; deploy on every merge)
       \     / \     / \     /
        ●───●   ●───●   ●───●   (feature branch → PR → CI → merge)
\end{verbatim}

Good for: web apps with continuous deployment, small/medium teams. Simple, fast. Assumes strong CI and feature flags for unfinished work.

\subsection{GitLab Flow (GitHub Flow + environment/release branches)}\label{gitlab-flow-github-flow--environmentrelease-branches}

Adds environment branches (\texttt{staging}, \texttt{production}) or release branches to GitHub Flow, so code promotes via merges: \texttt{main\ →\ staging\ →\ production}. Bridges "always deploy" simplicity with the reality of staged rollouts and regulated releases.

\setcodelang{}

\begin{verbatim}
main ──▶ staging ──▶ production    (changes flow "downstream" by merge)
                       │
              release/* (for products that ship versions)
\end{verbatim}

\subsection{Trunk-Based Development (TBD) --- what big-tech actually does at scale}\label{trunk-based-development-tbd--what-big-tech-actually-does-at-scale}

Everyone commits to a single \textbf{trunk} (\texttt{main}) many times a day, in tiny increments, behind \textbf{feature flags}. Branches are either nonexistent or live hours-to-a-day. CI runs on every commit; merges are continuous, so there\textquotesingle s no "big bang" integration.

\setcodelang{}

\begin{verbatim}
trunk ●─●─●─●─●─●─●─●─●─●─●─●─●  (hundreds of small commits/day, all behind flags)
       └ tiny short-lived branches (hours), merged via PR, deleted immediately
\end{verbatim}

\begin{itemize}
\tightlist
\item
  \textbf{Feature flags} decouple \emph{deploy} from \emph{release}: incomplete code ships to prod dark (flag off) and is toggled on when ready. This is why you can integrate constantly without shipping broken features.
\item
  \textbf{Why it scales:} continuous small merges = trivial conflicts (you\textquotesingle re always near HEAD), no long-lived divergence, no "merge hell." It\textquotesingle s the only model that works at Google/Meta scale with thousands of engineers.
\item
  Requires: fast, reliable CI; feature flags; good test coverage; sometimes release branches cut \emph{from} trunk for stabilization.
\end{itemize}

\subsection{Comparison}\label{comparison}

\begin{longtable}[]{@{}
  >{\raggedright\arraybackslash}p{(\columnwidth - 10\tabcolsep) * \real{0.0722}}
  >{\raggedright\arraybackslash}p{(\columnwidth - 10\tabcolsep) * \real{0.2102}}
  >{\raggedright\arraybackslash}p{(\columnwidth - 10\tabcolsep) * \real{0.1133}}
  >{\raggedright\arraybackslash}p{(\columnwidth - 10\tabcolsep) * \real{0.1209}}
  >{\raggedright\arraybackslash}p{(\columnwidth - 10\tabcolsep) * \real{0.2864}}
  >{\raggedright\arraybackslash}p{(\columnwidth - 10\tabcolsep) * \real{0.1970}}@{}}
\toprule\noalign{}
\rowcolor{acc!16}
\begin{minipage}[b]{\linewidth}\raggedright
Workflow
\end{minipage} & \begin{minipage}[b]{\linewidth}\raggedright
Long-lived branches
\end{minipage} & \begin{minipage}[b]{\linewidth}\raggedright
Branch lifetime
\end{minipage} & \begin{minipage}[b]{\linewidth}\raggedright
Integration freq
\end{minipage} & \begin{minipage}[b]{\linewidth}\raggedright
Best for
\end{minipage} & \begin{minipage}[b]{\linewidth}\raggedright
Main risk
\end{minipage} \\
\midrule\noalign{}
\endhead
\bottomrule\noalign{}
\endlastfoot
\textbf{Gitflow} & main + develop (+release/hotfix) & days--weeks & low & versioned/installed software, scheduled releases & merge hell, complexity \\
\rowcolor{zebra}
\textbf{GitHub Flow} & main & hours--days & high & continuous-deploy web apps & needs strong CI \\
\textbf{GitLab Flow} & main + env/release & hours--days & high & staged rollouts, compliance & more branches to manage \\
\rowcolor{zebra}
\textbf{Trunk-Based} & main only & hours & continuous & large teams, high deploy velocity & needs feature flags + great CI \\
\end{longtable}

\begin{quote}
\textbf{Interview takeaway:} "Small teams shipping continuously \(\rightarrow\) GitHub Flow / trunk-based with feature flags. Versioned products with release trains \(\rightarrow\) Gitflow. At FAANG scale, trunk-based with flags is the norm because continuous tiny merges avoid long-lived divergence and merge hell." Tie the choice to \emph{release cadence and CI maturity}, not personal taste.
\end{quote}

\section{Monorepo vs Polyrepo}\label{monorepo-vs-polyrepo}

\textbf{Monorepo:} all projects/services in \emph{one} repository (one history, one CI config, one set of tooling). \textbf{Polyrepo:} each project/service in its \emph{own} repository.

\setcodelang{}

\begin{verbatim}
Monorepo:                          Polyrepo:
  one-big-repo/                       repo: payments-service
    services/payments/                repo: auth-service
    services/auth/                    repo: web-frontend
    libs/common/                      repo: shared-lib  (versioned & published)
    web/frontend/                     ...each with its own CI, owners, releases
\end{verbatim}

\begin{longtable}[]{@{}
  >{\raggedright\arraybackslash}p{(\columnwidth - 4\tabcolsep) * \real{0.1671}}
  >{\raggedright\arraybackslash}p{(\columnwidth - 4\tabcolsep) * \real{0.4116}}
  >{\raggedright\arraybackslash}p{(\columnwidth - 4\tabcolsep) * \real{0.4214}}@{}}
\toprule\noalign{}
\rowcolor{acc!16}
\begin{minipage}[b]{\linewidth}\raggedright
Dimension
\end{minipage} & \begin{minipage}[b]{\linewidth}\raggedright
Monorepo
\end{minipage} & \begin{minipage}[b]{\linewidth}\raggedright
Polyrepo
\end{minipage} \\
\midrule\noalign{}
\endhead
\bottomrule\noalign{}
\endlastfoot
\textbf{Cross-project change} & \emph{Atomic} --- one commit changes a lib + all its callers, one PR, one review & Multi-repo dance: version-bump the lib, publish, update each consumer separately \\
\rowcolor{zebra}
\textbf{Shared tooling/CI} & One config, consistent everywhere & Each repo configures its own (drift) \\
\textbf{Dependency mgmt} & "Live at HEAD" --- everyone on the same version, no version hell & Semantic versioning, registries, diamond-dependency conflicts \\
\rowcolor{zebra}
\textbf{Code visibility / reuse} & Everything discoverable; easy refactors across boundaries & Siloed; reuse via published packages \\
\textbf{Build/scale} & Needs special tooling (Bazel, Buck) to not rebuild the world; clone gets huge & Each repo small and fast to clone/build \\
\rowcolor{zebra}
\textbf{Blast radius} & A bad commit/CI break can affect everyone & Naturally isolated per repo \\
\textbf{Access control} & Coarser (whole repo) --- needs path-based ACLs & Fine-grained per repo \\
\end{longtable}

\textbf{Why monorepos need heavy tooling:} a naive \texttt{git\ clone} of Google/Meta\textquotesingle s repo is impractical (hundreds of GB, millions of files). They rely on:

\begin{itemize}
\tightlist
\item
  \textbf{Sparse-checkout / partial clone} --- materialize only the directories you work in.
\item
  \textbf{Virtual filesystems} --- Meta\textquotesingle s \textbf{EdenFS}, Microsoft\textquotesingle s \textbf{VFS for Git / Scalar} --- files appear present but are fetched on access.
\item
  \textbf{Build systems} --- \textbf{Bazel} (Google), \textbf{Buck} (Meta) --- content-addressed, hermetic, incremental builds with remote caching so you don\textquotesingle t "rebuild the world."
\end{itemize}

\textbf{Real examples:}

\begin{itemize}
\tightlist
\item
  \textbf{Google:} a single monorepo (\texttt{google3}, billions of LOC) on Piper, "live at HEAD," one build system (Blaze/Bazel). Atomic cross-cutting refactors via large-scale change tooling.
\item
  \textbf{Meta:} monorepo on Sapling/EdenFS; sparse virtual checkouts; trunk-based.
\item
  \textbf{Microsoft:} Windows in Git via VFS for Git/Scalar (the famous 300 GB / 3.5M-file migration).
\item
  \textbf{Amazon:} the \emph{opposite} --- thousands of small service polyrepos, matching two-pizza team autonomy and independent deploy cadence.
\end{itemize}

\begin{quote}
\textbf{Interview takeaway:} The trade-off is \textbf{atomic cross-project changes + uniform tooling (monorepo)} vs \textbf{isolation + independent scaling + simpler per-repo ops (polyrepo)}. Monorepos buy consistency at the cost of needing serious build/VCS tooling; polyrepos buy isolation at the cost of dependency/version coordination. The "right" answer is "it depends on org structure (Conway\textquotesingle s Law) and tooling investment."
\end{quote}

\section{Hooks, Submodules, Subtrees, LFS, Sparse-Checkout, Signing}\label{hooks-submodules-subtrees-lfs-sparse-checkout-signing}

\subsection{Hooks}\label{hooks}

Scripts Git runs automatically at lifecycle points, stored in \texttt{.git/hooks/} (local, not committed) --- teams share them via tools like \textbf{pre-commit} or \textbf{Husky}.

\begin{longtable}[]{@{}
  >{\raggedright\arraybackslash}p{(\columnwidth - 4\tabcolsep) * \real{0.2379}}
  >{\raggedright\arraybackslash}p{(\columnwidth - 4\tabcolsep) * \real{0.3385}}
  >{\raggedright\arraybackslash}p{(\columnwidth - 4\tabcolsep) * \real{0.4236}}@{}}
\toprule\noalign{}
\rowcolor{acc!16}
\begin{minipage}[b]{\linewidth}\raggedright
Hook
\end{minipage} & \begin{minipage}[b]{\linewidth}\raggedright
Fires
\end{minipage} & \begin{minipage}[b]{\linewidth}\raggedright
Common use
\end{minipage} \\
\midrule\noalign{}
\endhead
\bottomrule\noalign{}
\endlastfoot
\texttt{pre-commit} & before a commit is created & lint, format, run fast tests, block secrets \\
\rowcolor{zebra}
\texttt{commit-msg} & after message entered & enforce conventional-commit format \\
\texttt{pre-push} & before \texttt{push} & run test suite, block pushing WIP/\texttt{-\/-fixup!} \\
\rowcolor{zebra}
\texttt{post-merge} / \texttt{post-checkout} & after merge/checkout & re-install deps, rebuild \\
server-side \texttt{pre-receive} & on the server before accepting a push & enforce policy (signed commits, no force-push to main) \\
\end{longtable}

Because local hooks aren\textquotesingle t committed and can be skipped (\texttt{-\/-no-verify}), real enforcement lives in \textbf{CI} (server-side). Local hooks are for fast feedback; CI is the gate.

\subsection{Submodules vs Subtrees}\label{submodules-vs-subtrees}

Both embed one repo inside another, differently:

\begin{itemize}
\tightlist
\item
  \textbf{Submodule} --- a \emph{pointer} (a recorded commit hash) to an external repo, checked out into a subdirectory. The parent repo stores only the reference, not the code. Pros: clear separation, the child has its own history. Cons: notoriously fiddly --- \texttt{git\ clone\ -\/-recursive}, \texttt{git\ submodule\ update\ -\/-init}, detached HEADs, easy to forget to commit the pointer bump.
\end{itemize}

\setcodelang{Bash}

\begin{Shaded}
\begin{Highlighting}[]
\FunctionTok{git}\NormalTok{ submodule add https://github.com/org/lib vendor/lib}
\FunctionTok{git}\NormalTok{ submodule update }\AttributeTok{{-}{-}init} \AttributeTok{{-}{-}recursive}
\end{Highlighting}
\end{Shaded}

\begin{itemize}
\tightlist
\item
  \textbf{Subtree} --- the external repo\textquotesingle s \emph{files and history are merged into} a subdirectory of the parent (no extra metadata for cloners). Pros: cloners need nothing special; the code is just there. Cons: pulling upstream updates uses special subtree commands; history is intertwined.
\end{itemize}

\setcodelang{Bash}

\begin{Shaded}
\begin{Highlighting}[]
\FunctionTok{git}\NormalTok{ subtree add }\AttributeTok{{-}{-}prefix}\OperatorTok{=}\NormalTok{vendor/lib https://github.com/org/lib main }\AttributeTok{{-}{-}squash}
\FunctionTok{git}\NormalTok{ subtree pull }\AttributeTok{{-}{-}prefix}\OperatorTok{=}\NormalTok{vendor/lib https://github.com/org/lib main }\AttributeTok{{-}{-}squash}
\end{Highlighting}
\end{Shaded}

\subsection{Git LFS (Large File Storage)}\label{git-lfs-large-file-storage}

Git is terrible at large binaries (every version is stored whole; they don\textquotesingle t delta well and bloat every clone forever). \textbf{LFS} replaces big files (videos, models, PSDs) with tiny text \emph{pointers} in the repo and stores the actual bytes on a separate LFS server, fetched on checkout.

\setcodelang{Bash}

\begin{Shaded}
\begin{Highlighting}[]
\FunctionTok{git}\NormalTok{ lfs install}
\FunctionTok{git}\NormalTok{ lfs track }\StringTok{"*.psd"} \StringTok{"*.bin"}    \CommentTok{\# writes patterns to .gitattributes}
\FunctionTok{git}\NormalTok{ add .gitattributes}
\end{Highlighting}
\end{Shaded}

\subsection{Sparse-checkout}\label{sparse-checkout}

Materialize only part of a (large/mono) repo\textquotesingle s tree into your working dir, while history still references everything.

\setcodelang{Bash}

\begin{Shaded}
\begin{Highlighting}[]
\FunctionTok{git}\NormalTok{ clone }\AttributeTok{{-}{-}filter}\OperatorTok{=}\NormalTok{blob:none }\AttributeTok{{-}{-}sparse} \OperatorTok{\textless{}}\NormalTok{url}\OperatorTok{\textgreater{}}   \CommentTok{\# partial clone: skip blobs until needed}
\FunctionTok{git}\NormalTok{ sparse{-}checkout set services/payments libs/common}
\end{Highlighting}
\end{Shaded}

\subsection{Signing (provenance / supply-chain trust)}\label{signing-provenance--supply-chain-trust}

Hashes prove \emph{integrity} (content wasn\textquotesingle t corrupted), not \emph{authenticity} (who made it). \textbf{Signed commits/tags} prove authorship.

\setcodelang{Bash}

\begin{Shaded}
\begin{Highlighting}[]
\FunctionTok{git}\NormalTok{ config commit.gpgsign true}
\FunctionTok{git}\NormalTok{ commit }\AttributeTok{{-}S} \AttributeTok{{-}m} \StringTok{"Release 2.4.0"}          \CommentTok{\# GPG{-}sign}
\FunctionTok{git}\NormalTok{ tag }\AttributeTok{{-}s}\NormalTok{ v2.4.0 }\AttributeTok{{-}m} \StringTok{"v2.4.0"}             \CommentTok{\# signed annotated tag}
\FunctionTok{git}\NormalTok{ log }\AttributeTok{{-}{-}show{-}signature}
\CommentTok{\# Modern alternative: SSH{-}key signing (no GPG needed)}
\FunctionTok{git}\NormalTok{ config gpg.format ssh}
\FunctionTok{git}\NormalTok{ config user.signingkey \textasciitilde{}/.ssh/id\_ed25519.pub}
\end{Highlighting}
\end{Shaded}

GitHub shows a "Verified" badge for signed commits --- increasingly required for release tags and supply-chain security (e.g., Sigstore/gitsign).

\section{Disasters \& Recovery (Worked Scenarios)}\label{disasters--recovery-worked-scenarios}

Stay calm; almost everything committed is recoverable via \textbf{reflog}. Here are the seven situations interviewers love.

\subsection{\texorpdfstring{1) I committed to \texttt{main} instead of my feature branch}{1) I committed to main instead of my feature branch}}\label{1-i-committed-to-main-instead-of-my-feature-branch}

Your last 2 commits should have been on a branch. \texttt{main} hasn\textquotesingle t been pushed yet.

\setcodelang{Bash}

\begin{Shaded}
\begin{Highlighting}[]
\FunctionTok{git}\NormalTok{ branch feature/login            }\CommentTok{\# create feature branch pointing at current (correct) state}
\FunctionTok{git}\NormalTok{ reset }\AttributeTok{{-}{-}hard}\NormalTok{ origin/main        }\CommentTok{\# rewind LOCAL main back to the remote (drops the 2 commits FROM main)}
\FunctionTok{git}\NormalTok{ switch feature/login            }\CommentTok{\# your commits live safely here now}
\end{Highlighting}
\end{Shaded}

If you\textquotesingle d already pushed those commits to \texttt{main} on a shared remote, \textbf{do not} reset --- \texttt{cherry-pick} them onto the branch and \texttt{revert} them on main instead (see \#2).

\subsection{2) Undo a commit already pushed to a shared branch}\label{2-undo-a-commit-already-pushed-to-a-shared-branch}

History is shared --- \textbf{never} \texttt{reset}/force-push. Use \texttt{revert} (forward-moving inverse commit):

\setcodelang{Bash}

\begin{Shaded}
\begin{Highlighting}[]
\FunctionTok{git}\NormalTok{ revert }\OperatorTok{\textless{}}\NormalTok{bad{-}sha}\OperatorTok{\textgreater{}}      \CommentTok{\# creates an inverse commit}
\FunctionTok{git}\NormalTok{ push origin main      }\CommentTok{\# safe, fast{-}forward push; no rewriting, no clobbering teammates}
\CommentTok{\# Reverting a merge commit:}
\FunctionTok{git}\NormalTok{ revert }\AttributeTok{{-}m}\NormalTok{ 1 }\OperatorTok{\textless{}}\NormalTok{merge{-}sha}\OperatorTok{\textgreater{}}   \CommentTok{\# {-}m 1 keeps the mainline parent}
\end{Highlighting}
\end{Shaded}

\subsection{\texorpdfstring{3) Accidental \texttt{reset\ -\/-hard}, lost work}{3) Accidental reset -\/-hard, lost work}}\label{3-accidental-reset---hard-lost-work}

The commits are unreachable but still in the object store (until GC).

\setcodelang{Bash}

\begin{Shaded}
\begin{Highlighting}[]
\FunctionTok{git}\NormalTok{ reflog                        }\CommentTok{\# find the hash from BEFORE the reset, e.g. a1b2c3d HEAD@\{1\}}
\FunctionTok{git}\NormalTok{ reset }\AttributeTok{{-}{-}hard}\NormalTok{ a1b2c3d          }\CommentTok{\# restore exactly}
\CommentTok{\# Or, if you only need it on a new branch:}
\FunctionTok{git}\NormalTok{ switch }\AttributeTok{{-}c}\NormalTok{ recovered a1b2c3d}
\end{Highlighting}
\end{Shaded}

Caveat: if \texttt{-\/-hard} wiped \emph{uncommitted} edits (never committed), reflog can\textquotesingle t help --- those bytes were never an object.

\subsection{4) Force-pushed and clobbered a teammate\textquotesingle s commits}\label{4-force-pushed-and-clobbered-a-teammates-commits}

You \texttt{git\ push\ -\/-force} and overwrote commits the teammate had pushed. The remote\textquotesingle s old tip is gone \emph{server-side}, but it likely still exists in \emph{someone\textquotesingle s} local reflog or \texttt{origin/...} tracking ref.

\setcodelang{Bash}

\begin{Shaded}
\begin{Highlighting}[]
\CommentTok{\# Teammate (or anyone who fetched the old state) finds the old tip:}
\FunctionTok{git}\NormalTok{ reflog show origin/main          }\CommentTok{\# or check their local branch reflog}
\FunctionTok{git}\NormalTok{ push }\AttributeTok{{-}{-}force{-}with{-}lease}\NormalTok{ origin }\OperatorTok{\textless{}}\NormalTok{old{-}good{-}sha}\OperatorTok{\textgreater{}}\NormalTok{:main   }\CommentTok{\# restore}
\end{Highlighting}
\end{Shaded}

\textbf{Prevention:} always use \textbf{\texttt{git\ push\ -\/-force-with-lease}} instead of \texttt{-\/-force}. \texttt{-\/-force-with-lease} refuses the push if the remote moved since your last fetch (i.e., someone else pushed) --- it \emph{won\textquotesingle t} silently clobber. And protect \texttt{main}/release branches server-side (GitHub branch protection: no force-push, require PRs).

\subsection{5) Committed a secret / huge file --- purge from ALL history}\label{5-committed-a-secret--huge-file--purge-from-all-history}

Removing it in a new commit is \textbf{not enough} --- it\textquotesingle s still in history (and a leaked secret must be considered compromised \(\rightarrow\) \textbf{rotate the credential immediately}, then scrub). Rewriting history requires \texttt{git\ filter-repo} (recommended) or \textbf{BFG}:

\setcodelang{Bash}

\begin{Shaded}
\begin{Highlighting}[]
\CommentTok{\# Preferred: git{-}filter{-}repo (faster, safer than the deprecated filter{-}branch)}
\FunctionTok{git}\NormalTok{ filter{-}repo }\AttributeTok{{-}{-}path}\NormalTok{ config/secrets.yml }\AttributeTok{{-}{-}invert{-}paths}   \CommentTok{\# strip a file from ALL history}
\FunctionTok{git}\NormalTok{ filter{-}repo }\AttributeTok{{-}{-}strip{-}blobs{-}bigger{-}than}\NormalTok{ 50M              }\CommentTok{\# strip oversized blobs}

\CommentTok{\# Alternative: BFG Repo{-}Cleaner (simpler for the common cases)}
\ExtensionTok{bfg} \AttributeTok{{-}{-}delete{-}files}\NormalTok{ secrets.yml}
\ExtensionTok{bfg} \AttributeTok{{-}{-}strip{-}blobs{-}bigger{-}than}\NormalTok{ 50M}
\FunctionTok{git}\NormalTok{ reflog expire }\AttributeTok{{-}{-}expire}\OperatorTok{=}\NormalTok{now }\AttributeTok{{-}{-}all} \KeywordTok{\&\&} \FunctionTok{git}\NormalTok{ gc }\AttributeTok{{-}{-}prune}\OperatorTok{=}\NormalTok{now }\AttributeTok{{-}{-}aggressive}
\end{Highlighting}
\end{Shaded}

This \textbf{rewrites every commit hash} \(\rightarrow\) everyone must re-clone (a coordinated, disruptive operation). Then \textbf{force-push} (with team coordination) and \textbf{invalidate the leaked secret}. Order of operations: rotate the secret first, \emph{then} scrub history, \emph{then} notify the team to re-clone.

\subsection{6) Detached HEAD --- I made commits and now they\textquotesingle re "gone"}\label{6-detached-head--i-made-commits-and-now-theyre-gone}

You \texttt{git\ checkout\ \textless{}sha\textgreater{}} (detached), committed, then switched away. Those commits aren\textquotesingle t on any branch.

\setcodelang{Bash}

\begin{Shaded}
\begin{Highlighting}[]
\CommentTok{\# Right after committing in detached state, BEFORE switching:}
\FunctionTok{git}\NormalTok{ switch }\AttributeTok{{-}c}\NormalTok{ rescue        }\CommentTok{\# attach a branch to current HEAD → commits saved}

\CommentTok{\# If you ALREADY switched away:}
\FunctionTok{git}\NormalTok{ reflog                  }\CommentTok{\# find the dangling commit you made}
\FunctionTok{git}\NormalTok{ switch }\AttributeTok{{-}c}\NormalTok{ rescue }\OperatorTok{\textless{}}\NormalTok{that{-}sha}\OperatorTok{\textgreater{}}
\end{Highlighting}
\end{Shaded}

A detached HEAD itself is harmless (Git warns you); the danger is only \emph{losing track} of commits you make in it.

\subsection{7) A gnarly conflict during a big merge/rebase}\label{7-a-gnarly-conflict-during-a-big-mergerebase}

\setcodelang{Bash}

\begin{Shaded}
\begin{Highlighting}[]
\FunctionTok{git}\NormalTok{ config merge.conflictstyle zdiff3   }\CommentTok{\# show the base too — see who changed what}
\FunctionTok{git}\NormalTok{ config rerere.enabled true          }\CommentTok{\# remember resolutions for repeats}
\FunctionTok{git}\NormalTok{ status                              }\CommentTok{\# list conflicted files}
\FunctionTok{git}\NormalTok{ mergetool                           }\CommentTok{\# launch a 3{-}way visual tool (vimdiff, meld, vscode...)}
\CommentTok{\# Resolve each file; remove ALL \textless{}\textless{}\textless{}\textless{} |||| ==== \textgreater{}\textgreater{}\textgreater{}\textgreater{} markers; then:}
\FunctionTok{git}\NormalTok{ add }\OperatorTok{\textless{}}\NormalTok{file}\OperatorTok{\textgreater{}}
\FunctionTok{git}\NormalTok{ merge }\AttributeTok{{-}{-}continue}      \CommentTok{\# or: git rebase {-}{-}continue}
\CommentTok{\# Escape hatches:}
\FunctionTok{git}\NormalTok{ merge }\AttributeTok{{-}{-}abort}
\FunctionTok{git}\NormalTok{ rebase }\AttributeTok{{-}{-}abort}
\CommentTok{\# Per{-}file "just take one side":}
\FunctionTok{git}\NormalTok{ checkout }\AttributeTok{{-}{-}ours}\NormalTok{  config.py }\KeywordTok{\&\&} \FunctionTok{git}\NormalTok{ add config.py}
\FunctionTok{git}\NormalTok{ checkout }\AttributeTok{{-}{-}theirs}\NormalTok{ config.py }\KeywordTok{\&\&} \FunctionTok{git}\NormalTok{ add config.py}
\end{Highlighting}
\end{Shaded}

For a \emph{huge} conflicted rebase, consider \texttt{git\ rebase\ -\/-onto} to replay only the commits you need, or split it into smaller steps. If a conflict recurs every rebase, \texttt{rerere} saves your sanity.

\begin{quote}
\textbf{Interview takeaway:} The meta-answer to "you broke the repo, what now?" is: \textbf{don\textquotesingle t panic, \texttt{git\ reflog} first, prefer non-destructive fixes (\texttt{revert}, \texttt{-\/-force-with-lease}), and never rewrite shared history without team coordination.} That composure is exactly what they\textquotesingle re testing.
\end{quote}

\section{Architecture / Diagrams}\label{architecture--diagrams}

\subsection{The full data flow}\label{the-full-data-flow}

\setcodelang{}

\begin{verbatim}
   ┌────────────────────────────────────────────────────────────────────┐
   │                         YOUR LOCAL MACHINE                          │
   │                                                                     │
   │  Working Dir ──add──▶ Index ──commit──▶ Local Repo (.git/objects)   │
   │   (files)            (staging)          ├─ blobs / trees / commits  │
   │      ▲                                  ├─ refs/heads/* (branches)  │
   │      └────── restore / checkout ────────┤  HEAD, reflog             │
   │                                         └─ refs/remotes/origin/*    │
   └───────────────────────────────────┬─────────────────▲──────────────┘
                                 push   │                 │  fetch / pull
                                        ▼                 │
   ┌────────────────────────────────────────────────────────────────────┐
   │                    REMOTE (GitHub / GitLab / origin)                │
   │   refs/heads/*  ◀── PRs, code review, branch protection, CI/CD ──▶  │
   └────────────────────────────────────────────────────────────────────┘
\end{verbatim}

\subsection{Object graph for a small history}\label{object-graph-for-a-small-history}

\setcodelang{}

\begin{verbatim}
   tags/v1.0 ──┐
               ▼
  C1 ◀── C2 ◀── C3 ◀── C4   ← refs/heads/main ◀── HEAD
                  ▲
                  └── C3a ◀── C3b   ← refs/heads/feature

  Each Cn ──▶ a tree ──▶ blobs (unchanged files share blobs across commits)
  Arrows point to PARENTS (commits know their past, not their future).
\end{verbatim}

\subsection{Merge vs Rebase topology (recap)}\label{merge-vs-rebase-topology-recap}

\setcodelang{}

\begin{verbatim}
MERGE:                              REBASE:
   A─B─┬─C─D──┐                        A─B─E─F─C'─D'
       │      M (2 parents)            (linear; C',D' are NEW commits)
       └─E─F──┘
\end{verbatim}

\section{Real-World Examples}\label{real-world-examples}

\begin{itemize}
\tightlist
\item
  \textbf{Hotfix backport (cherry-pick):} A null-pointer crash is fixed on \texttt{main} (commit \texttt{a1b2c3}). Customers on the \texttt{release/3.2} line need it now. \texttt{git\ switch\ release/3.2\ \&\&\ git\ cherry-pick\ a1b2c3} \(\rightarrow\) tag \texttt{v3.2.1} \(\rightarrow\) ship. The full \texttt{main} (with unrelated risky features) never touches the release branch.
\item
  \textbf{Finding a perf regression (bisect):} A dashboard got 4\(\times\) slower "sometime last sprint." \texttt{git\ bisect\ start;\ git\ bisect\ bad\ HEAD;\ git\ bisect\ good\ v2.7.0;\ git\ bisect\ run\ ./bench.sh} walks \textasciitilde9 commits and prints the offending commit --- a seemingly innocent change to a query that dropped an index hint.
\item
  \textbf{Cleaning a PR (interactive rebase):} 14 "WIP", "fix", "oops" commits become 3 logical commits via \texttt{git\ rebase\ -i} (squash/fixup/reword) before requesting review. Reviewers see intent, not your keystroke history.
\item
  \textbf{Leaked AWS key (filter-repo):} An engineer pushes \texttt{.env} with a live key. Response: rotate the key (assume compromised), \texttt{git\ filter-repo\ -\/-path\ .env\ -\/-invert-paths}, force-push with team coordination, everyone re-clones, add \texttt{.env} to \texttt{.gitignore} and a \texttt{pre-commit} secret scanner.
\item
  \textbf{Microsoft Windows on Git:} \textasciitilde300 GB, 3.5M files, 4,000+ engineers --- impossible with vanilla Git; solved by inventing VFS for Git (now Scalar) with on-demand file hydration and partial clone.
\item
  \textbf{Google "live at HEAD":} one monorepo, everyone builds against trunk; a library change and all its callers update in a \emph{single atomic commit} --- impossible in a polyrepo without coordinated multi-repo releases.
\end{itemize}

\section{Real-Life Analogies}\label{real-life-analogies}

\begin{itemize}
\tightlist
\item
  \textbf{Commits = save points in a video game.} Each is a full snapshot you can reload. The hash is the save-slot ID; the parent link is "the save you made before this one."
\item
  \textbf{Branches = bookmarks / sticky notes} on the timeline. Moving a branch = peeling the sticky note off one commit and putting it on another. Cheap because it\textquotesingle s just a sticky note, not a copy of the book.
\item
  \textbf{Staging area = a shopping cart.} Your working dir is the whole store (all your changes). You put \emph{selected} items in the cart (\texttt{git\ add}), then check out (\texttt{git\ commit}). \texttt{git\ add\ -p} is picking individual items off a shelf rather than sweeping the whole shelf in.
\item
  \textbf{Merge = combining two edited copies of a document} by hand, with the \emph{original} as reference (3-way) so you know who changed what. A conflict is "you both rewrote the same paragraph differently."
\item
  \textbf{Rebase = re-recording your podcast episodes on top of the newest intro} --- same content, brand-new recordings (new hashes). Don\textquotesingle t re-record episodes listeners already downloaded (the Golden Rule).
\item
  \textbf{Reflog = your browser history for HEAD.} Even after you "close the tab" (delete a branch / hard reset), the history remembers where you were, so you can go back.
\item
  \textbf{Cherry-pick = copying one paragraph from one document into another}, rather than merging the whole document.
\item
  \textbf{\texttt{revert} = an erratum/correction note} appended at the end ("the change on page 5 was wrong, here\textquotesingle s the fix") vs \textbf{\texttt{reset} = ripping the page out} as if it never existed.
\item
  \textbf{Content addressing = a coat check that gives you a ticket equal to a fingerprint of the coat} --- identical coats get identical tickets, and you can\textquotesingle t swap a coat without the ticket changing.
\end{itemize}

\section{Memory Tricks / Mnemonics}\label{memory-tricks--mnemonics}

\begin{itemize}
\tightlist
\item
  \textbf{Four objects: "BTCT" --- Blob, Tree, Commit, Tag.} (Blob = bytes, Tree = directory, Commit = snapshot+parent, Tag = name.)
\item
  \textbf{Three areas: "Work, Stage, Store"} (working dir \(\rightarrow\) index \(\rightarrow\) repo). Verbs: \textbf{add} moves right, \textbf{restore} moves left.
\item
  \textbf{Reset modes --- "Soft keeps Staged, Mixed keeps Modified, Hard kills it":}

  \begin{itemize}
  \tightlist
  \item
    \textbf{S}oft \(\rightarrow\) index (\textbf{S}taged).
  \item
    \textbf{M}ixed \(\rightarrow\) working dir (\textbf{M}odified/unstaged).
  \item
    \textbf{H}ard \(\rightarrow\) gone.
  \end{itemize}
\item
  \textbf{Reset vs Revert: "Reset Rewinds (private), Revert Records (public)."}
\item
  \textbf{Golden Rule: "Never rebase what you\textquotesingle ve pushed."} (Or: \emph{"Rebase local, merge shared."})
\item
  \textbf{Pull default: "pull = fetch + merge; add \texttt{-\/-rebase} for a clean line."}
\item
  \textbf{Disaster reflex: "Reflog first, panic never."}
\item
  \textbf{Pick the tool: "Regression\(\rightarrow\)Bisect, Lost\(\rightarrow\)Reflog, Backport\(\rightarrow\)Cherry-pick, Pause\(\rightarrow\)Stash."}
\item
  \textbf{\texttt{-\/-force-with-lease}, not \texttt{-\/-force}:} "lease checks the lock before you smash it."
\end{itemize}

\section{Common Interview Questions}\label{common-interview-questions}

\subsection{Q1: Does Git store diffs or snapshots? Why does it matter?}\label{q1-does-git-store-diffs-or-snapshots-why-does-it-matter}

\textbf{Model answer:} Git stores \textbf{full snapshots} per commit, with unchanged files pointing to the same blob (so it\textquotesingle s not wasteful). Conceptually each commit is a complete tree; physically, packfiles add delta compression. This makes branch/checkout/diff between \emph{any} two commits fast and uniform, and it\textquotesingle s why merging is cheaper than in delta-based systems where you\textquotesingle d replay a chain of changes. Most older systems (SVN/CVS) store a base plus per-file deltas, so reconstructing an old state means replaying changes.

\textbf{Follow-ups:}

\begin{itemize}
\tightlist
\item
  \emph{How does it avoid storing duplicate content?} \(\rightarrow\) Content addressing: identical bytes hash to the same blob, stored once; packfiles further delta-compress similar objects.
\item
  \emph{Then what are packfiles?} \(\rightarrow\) A storage optimization (\texttt{git\ gc}) that packs loose objects and delta-compresses similar ones; the logical model stays "snapshots."
\end{itemize}

\subsection{Q2: Walk me through the four Git object types and how a commit hash is formed.}\label{q2-walk-me-through-the-four-git-object-types-and-how-a-commit-hash-is-formed}

\textbf{Model answer:} \textbf{Blob} = file contents (no name). \textbf{Tree} = a directory listing mapping names+modes to blobs/subtrees. \textbf{Commit} = a pointer to one root tree + parent commit(s) + author/committer/message. \textbf{Tag} = a named (optionally signed) pointer. A commit\textquotesingle s hash is \texttt{SHA-1("commit\ \textless{}len\textgreater{}\textbackslash{}0"\ +\ commit-text)}, and the commit text \emph{includes the tree hash and the parent hash}. Because the parent\textquotesingle s hash is baked in, changing any ancestor changes that commit\textquotesingle s hash and cascades to all descendants --- that\textquotesingle s the tamper-evidence and why the same diff at a different parent yields a different commit.

\textbf{Follow-ups:}

\begin{itemize}
\tightlist
\item
  \emph{Where does a filename live?} \(\rightarrow\) In the tree entry, not the blob --- so a pure rename reuses the blob.
\item
  \emph{SHA-1 is broken --- is Git insecure?} \(\rightarrow\) Git uses collision-detecting SHA-1 and supports SHA-256 repos; the hash is for integrity/dedup, not an authenticity guarantee --- that\textquotesingle s what signing is for.
\end{itemize}

\subsection{Q3: Explain rebase vs merge and when you\textquotesingle d use each.}\label{q3-explain-rebase-vs-merge-and-when-youd-use-each}

\textbf{Model answer:} Merge integrates by creating a \textbf{merge commit} with two parents --- it preserves true history and never rewrites commits. Rebase \textbf{replays} your commits onto a new base, creating \emph{new commits with new hashes} for a \textbf{linear} history. I rebase my \emph{local} feature branch to keep it clean and easy to review/bisect, and to incorporate the latest \texttt{main} without a merge bubble. I \textbf{merge} to integrate shared branches and to preserve "these commits were one feature." The hard rule: never rebase commits already pushed/shared, because rewriting their hashes breaks anyone who built on the originals.

\textbf{Follow-ups:}

\begin{itemize}
\tightlist
\item
  \emph{Why exactly does rebasing shared history break people?} \(\rightarrow\) It replaces \texttt{C,D} with \texttt{C\textquotesingle{},D\textquotesingle{}}; teammates who based work on \texttt{C,D} now diverge --- duplicate commits, messy merges, and a force-push can clobber work.
\item
  \emph{\texttt{git\ pull} default?} \(\rightarrow\) fetch+merge; I prefer \texttt{pull\ -\/-rebase} to avoid noise merge commits on local-only work.
\end{itemize}

\subsection{\texorpdfstring{Q4: \texttt{reset\ -\/-soft} vs \texttt{-\/-mixed} vs \texttt{-\/-hard}?}{Q4: reset -\/-soft vs -\/-mixed vs -\/-hard?}}\label{q4-reset---soft-vs---mixed-vs---hard}

\textbf{Model answer:} All move the branch pointer to a target commit; they differ in how far they reset. \texttt{-\/-soft} moves HEAD only, leaving the undone changes \textbf{staged} (good for re-doing/combining the last N commits). \texttt{-\/-mixed} (default) also resets the index, leaving changes as \textbf{unstaged} edits. \texttt{-\/-hard} also resets the working dir --- \textbf{destroying} uncommitted changes. \texttt{-\/-soft}/\texttt{-\/-mixed} are safe; \texttt{-\/-hard} can lose uncommitted work permanently (reflog recovers commits, not never-committed edits).

\textbf{Follow-ups:}

\begin{itemize}
\tightlist
\item
  \emph{Undo the last 3 commits but keep all the code to recommit as one?} \(\rightarrow\) \texttt{git\ reset\ -\/-soft\ HEAD\textasciitilde{}3} then commit.
\item
  \emph{You hard-reset and lost commits --- recover?} \(\rightarrow\) \texttt{git\ reflog} \(\rightarrow\) \texttt{git\ reset\ -\/-hard\ \textless{}pre-reset-sha\textgreater{}}.
\end{itemize}

\subsection{Q5: How do you undo a commit that\textquotesingle s already been pushed to a shared branch?}\label{q5-how-do-you-undo-a-commit-thats-already-been-pushed-to-a-shared-branch}

\textbf{Model answer:} Use \texttt{git\ revert\ \textless{}sha\textgreater{}}, which creates a \emph{new} commit applying the inverse diff, then push normally. It moves history forward and doesn\textquotesingle t rewrite shared commits, so no one is disrupted and no force-push is needed. \texttt{reset} + force-push would rewrite shared history and break collaborators. For a merge commit, \texttt{git\ revert\ -m\ 1\ \textless{}merge-sha\textgreater{}} to keep the mainline parent.

\textbf{Follow-ups:}

\begin{itemize}
\tightlist
\item
  \emph{Later you want the change back?} \(\rightarrow\) Revert the revert, or cherry-pick the original.
\item
  \emph{When is \texttt{reset} acceptable then?} \(\rightarrow\) Only on local, unpushed commits nobody depends on.
\end{itemize}

\subsection{Q6: A bug appeared "sometime in the last 200 commits." How do you find it fast?}\label{q6-a-bug-appeared-sometime-in-the-last-200-commits-how-do-you-find-it-fast}

\textbf{Model answer:} \texttt{git\ bisect} --- a binary search over the commit range. Mark a known-good commit and the current bad one; Git checks out the midpoint, I test it and mark good/bad, and it halves the range each step --- \textasciitilde8 tests for 200 commits instead of 200. I\textquotesingle d automate it with \texttt{git\ bisect\ run\ \textless{}test-script\textgreater{}} (exit 0 = good, non-zero = bad, 125 = skip) so Git finds the first bad commit unattended. The output is the exact offending commit and its diff.

\textbf{Follow-ups:}

\begin{itemize}
\tightlist
\item
  \emph{What makes bisect work well?} \(\rightarrow\) Small, individually-building commits; if a commit doesn\textquotesingle t build, mark it \texttt{skip} (exit 125).
\item
  \emph{Complexity?} \(\rightarrow\) O(log n) tests.
\end{itemize}

\subsection{\texorpdfstring{Q7: How do you take a single fix from \texttt{main} onto a release branch?}{Q7: How do you take a single fix from main onto a release branch?}}\label{q7-how-do-you-take-a-single-fix-from-main-onto-a-release-branch}

\textbf{Model answer:} \texttt{git\ cherry-pick\ \textless{}fix-sha\textgreater{}} while on the release branch --- it applies just that commit\textquotesingle s diff as a new commit, without dragging in unrelated changes. I\textquotesingle d use \texttt{-x} to record the original SHA in the message for traceability, then tag a patch release. It\textquotesingle s the standard backport mechanism. Caveat: it creates a duplicate commit (different hash), so I avoid using cherry-pick as a substitute for actually merging long-running branches.

\textbf{Follow-ups:}

\begin{itemize}
\tightlist
\item
  \emph{Range of commits?} \(\rightarrow\) \texttt{git\ cherry-pick\ A\^{}..B}.
\item
  \emph{Conflict during pick?} \(\rightarrow\) resolve, \texttt{git\ add}, \texttt{git\ cherry-pick\ -\/-continue} (or \texttt{-\/-abort}).
\end{itemize}

\subsection{Q8: Explain the 3-way merge and why conflicts happen.}\label{q8-explain-the-3-way-merge-and-why-conflicts-happen}

\textbf{Model answer:} Git finds the \textbf{merge base} (nearest common ancestor) and compares three versions of each file: base, ours, theirs. If only one side changed a region, Git takes that change automatically; if \emph{both} sides changed the \emph{same} region relative to the base, Git can\textquotesingle t safely choose, so it marks a \textbf{conflict} with \texttt{\textless{}\textless{}\textless{}\textless{}\ ====\ \textgreater{}\textgreater{}\textgreater{}\textgreater{}} markers and asks me to resolve. The base is what makes auto-merge possible --- a plain two-way diff can\textquotesingle t tell which side made a change. The default strategy is \textbf{ort} (a faster, more correct successor to recursive), which handles multiple common ancestors by recursively merging them into one virtual base.

\textbf{Follow-ups:}

\begin{itemize}
\tightlist
\item
  \emph{Make conflicts easier?} \(\rightarrow\) \texttt{merge.conflictstyle=zdiff3} (shows the base) and \texttt{rerere.enabled} (reuse past resolutions).
\item
  \emph{Take one whole side per file?} \(\rightarrow\) \texttt{git\ checkout\ -\/-ours/-\/-theirs\ \textless{}file\textgreater{}} then \texttt{git\ add}.
\end{itemize}

\subsection{Q9: You committed a secret (API key) and pushed it. What now?}\label{q9-you-committed-a-secret-api-key-and-pushed-it-what-now}

\textbf{Model answer:} First, \textbf{rotate the credential immediately} --- once pushed, treat it as compromised; scrubbing alone is not enough. Then remove it from \emph{all} history with \texttt{git\ filter-repo\ -\/-path\ \textless{}file\textgreater{}\ -\/-invert-paths} (or BFG), expire reflogs and GC, and force-push with team coordination since this rewrites every subsequent hash and everyone must re-clone. Finally, add the file to \texttt{.gitignore} and add a \texttt{pre-commit}/CI secret scanner so it can\textquotesingle t recur. Order matters: rotate \(\rightarrow\) scrub \(\rightarrow\) re-clone.

\textbf{Follow-ups:}

\begin{itemize}
\tightlist
\item
  \emph{Why isn\textquotesingle t a "delete the file" commit enough?} \(\rightarrow\) The secret still lives in history and in the remote; anyone can check out the old commit.
\item
  \emph{filter-repo vs filter-branch?} \(\rightarrow\) filter-branch is slow and deprecated/error-prone; filter-repo (or BFG) is the recommended tool.
\end{itemize}

\subsection{Q10: What does "a branch is just a pointer" mean, and what\textquotesingle s a detached HEAD?}\label{q10-what-does-a-branch-is-just-a-pointer-mean-and-whats-a-detached-head}

\textbf{Model answer:} A branch is a 41-byte file under \texttt{.git/refs/heads/} holding one commit hash --- creating a branch writes that file (O(1)). HEAD is normally a symbolic ref to the current branch; committing appends a commit and advances that branch pointer. A \textbf{detached HEAD} is when HEAD points \emph{directly} at a commit instead of a branch (e.g., after \texttt{git\ checkout\ \textless{}sha\textgreater{}}). Commits made there aren\textquotesingle t tracked by any branch, so they\textquotesingle re easy to lose --- the fix is to attach a branch (\texttt{git\ switch\ -c\ name}) before leaving, or recover via reflog afterward.

\textbf{Follow-ups:}

\begin{itemize}
\tightlist
\item
  \emph{When do you intentionally detach?} \(\rightarrow\) To inspect/build an old commit, or in CI checking out a specific SHA.
\item
  \emph{Recover commits made while detached?} \(\rightarrow\) \texttt{git\ reflog} \(\rightarrow\) \texttt{git\ switch\ -c\ rescue\ \textless{}sha\textgreater{}}.
\end{itemize}

\subsection{Q11: Compare monorepo and polyrepo. Which would you choose?}\label{q11-compare-monorepo-and-polyrepo-which-would-you-choose}

\textbf{Model answer:} A monorepo holds everything in one repo, enabling \textbf{atomic cross-project changes} (change a library and all callers in one reviewed commit) and uniform tooling/CI, with "live at HEAD" dependencies --- but it needs serious build/VCS tooling (Bazel/Buck, sparse-checkout, virtual filesystems) to scale and has a larger blast radius. A polyrepo isolates each service: small fast repos, independent deploys, fine-grained access --- at the cost of cross-repo dependency/version coordination and tooling drift. I\textquotesingle d choose based on org structure (Conway\textquotesingle s Law) and tooling maturity: large, tightly-coupled codebases with platform teams favor monorepo (Google/Meta); autonomous service teams with independent cadence favor polyrepo (Amazon).

\textbf{Follow-ups:}

\begin{itemize}
\tightlist
\item
  \emph{Biggest monorepo pain?} \(\rightarrow\) Clone/build scale \(\rightarrow\) solved with partial clone, sparse-checkout, VFS, remote build cache.
\item
  \emph{Biggest polyrepo pain?} \(\rightarrow\) Diamond dependencies and coordinating a breaking change across many repos.
\end{itemize}

\subsection{Q12: Which branching workflow, and why?}\label{q12-which-branching-workflow-and-why}

\textbf{Model answer:} It depends on release cadence and CI maturity. For a continuously-deployed web service with strong CI, \textbf{trunk-based development with feature flags} (or GitHub Flow) --- everyone commits small changes to \texttt{main} many times a day behind flags, so integration is continuous and you avoid long-lived divergence and merge hell; flags decouple deploy from release. For versioned/installed software with release trains (mobile, libraries), \textbf{Gitflow} --- \texttt{main}/\texttt{develop} plus \texttt{release/*} and \texttt{hotfix/*} branches give controlled stabilization. The FAANG-scale default is trunk-based with flags.

\textbf{Follow-ups:}

\begin{itemize}
\tightlist
\item
  \emph{Why does trunk-based scale to thousands of engineers?} \(\rightarrow\) Tiny, continuous merges keep everyone near HEAD \(\rightarrow\) trivial conflicts; no big-bang integration.
\item
  \emph{How do you ship unfinished features on trunk?} \(\rightarrow\) Feature flags (dark launch), plus tests guarding flagged paths.
\end{itemize}

\section{Senior-Level Discussion Points}\label{senior-level-discussion-points}

\begin{itemize}
\tightlist
\item
  \textbf{Deploy vs release decoupling.} Feature flags let you \emph{deploy} code continuously while \emph{releasing} features independently --- the foundation of trunk-based development and progressive rollouts (canary, ring deployments). Mention flag-debt cleanup as the cost.
\item
  \textbf{Commit hygiene as an engineering practice.} Atomic, single-purpose, well-described commits aren\textquotesingle t vanity --- they make \texttt{bisect} precise, \texttt{revert} safe, \texttt{blame} meaningful, and reviews fast. "Each commit should compile and pass tests" is a real bar at top shops.
\item
  \textbf{Conventional Commits + SemVer + automated releases.} \texttt{feat:}, \texttt{fix:}, \texttt{feat!:}/\texttt{BREAKING\ CHANGE:} map mechanically to \textbf{Semantic Versioning} (MAJOR.MINOR.PATCH): \texttt{fix→patch}, \texttt{feat→minor}, breaking\(\rightarrow\)major. Tools (semantic-release, changesets) auto-bump versions and generate changelogs from commit messages. Shows you think about release automation, not just code.
\item
  \textbf{Force-push safety culture.} \texttt{-\/-force-with-lease} over \texttt{-\/-force}; branch protection (no force-push to main, required PRs, required green CI, required reviews, signed commits) --- these are guardrails that prevent the worst disasters at scale.
\item
  \textbf{CI as the real gate, hooks as fast feedback.} Local hooks are bypassable (\texttt{-\/-no-verify}) and not shared; server-side checks (CI, \texttt{pre-receive}) are the enforcement boundary. Articulating this distinction is a maturity signal.
\item
  \textbf{Stacked diffs / stacked PRs.} Breaking a big change into a \emph{stack} of dependent small PRs (Graphite, Phabricator, Sapling) keeps each review small while preserving logical sequence --- common at Meta and increasingly elsewhere. Rebasing the stack on \texttt{main} is routine.
\item
  \textbf{Merge queues.} At high commit volume, "green on my branch" isn\textquotesingle t enough because main moved; \textbf{merge queues} (GitHub merge queue, Bors, Zuul) re-test each PR against the \emph{latest} main in order before merging, preventing semantic conflicts and broken-main races.
\item
  \textbf{Supply-chain provenance.} Signed commits/tags (GPG/SSH/Sigstore-gitsign), protected tags, and reproducible builds tie artifacts to source --- increasingly mandatory (SLSA framework).
\item
  \textbf{Monorepo build economics.} The real cost of a monorepo isn\textquotesingle t Git, it\textquotesingle s \emph{builds} --- content-addressed, hermetic build systems with remote caching (Bazel/Buck) are what make "don\textquotesingle t rebuild the world" possible.
\end{itemize}

\section{Typical Mistakes Candidates Make}\label{typical-mistakes-candidates-make}

\begin{itemize}
\tightlist
\item
  \textbf{Treating Git as memorized incantations.} Not knowing \emph{why} \texttt{reset\ -\/-hard} is dangerous or what a branch physically is. Interviewers probe the model, not flags.
\item
  \textbf{Confusing \texttt{reset} and \texttt{revert}} --- proposing \texttt{reset} + force-push to undo something on a \emph{shared} branch. That breaks teammates; \texttt{revert} is the safe public undo.
\item
  \textbf{Not knowing what \texttt{-\/-hard} destroys} --- claiming reflog can recover \emph{uncommitted} edits wiped by \texttt{reset\ -\/-hard}/\texttt{clean}. It can\textquotesingle t; those were never objects.
\item
  \textbf{Saying rebase and merge are interchangeable.} They differ in topology and whether commits are rewritten --- and rebasing shared history is a cardinal sin.
\item
  \textbf{Forgetting the Golden Rule} or being unable to explain \emph{why} rebasing pushed commits breaks collaborators (the new-hash divergence).
\item
  \textbf{Thinking deleting a file in a new commit removes a leaked secret.} It\textquotesingle s still in history; you must rewrite history \emph{and} rotate the credential.
\item
  \textbf{Using \texttt{git\ push\ -\/-force}} instead of \texttt{-\/-force-with-lease}, and not mentioning branch protection.
\item
  \textbf{\texttt{git\ add\ .} everything, vague messages} ("fixes", "stuff") --- no atomicity, useless for blame/bisect/review.
\item
  \textbf{Resolving a conflict by deleting one side blindly} or leaving stray \texttt{\textless{}\textless{}\textless{}\textless{}}/\texttt{\textgreater{}\textgreater{}\textgreater{}\textgreater{}} markers in committed code.
\item
  \textbf{Not knowing \texttt{reflog} exists} --- panicking and re-cloning when recovery was one command away.
\item
  \textbf{Overusing cherry-pick} as a merge substitute, creating duplicate commits and later conflicts.
\item
  \textbf{Choosing a workflow dogmatically} rather than tying it to release cadence, team size, and CI maturity.
\end{itemize}

\section{How This Connects to Other Topics}\label{how-this-connects-to-other-topics}

\begin{itemize}
\tightlist
\item
  \textbf{Operating Systems:} Git is a content-addressable \emph{filesystem}; blobs/trees mirror inodes/directories, and packfiles are storage compaction. \texttt{.git/index} is a cache like a page table. File-watching tools (and VFS-for-Git) hook OS filesystem layers.
\item
  \textbf{Hashing \& Data Structures:} Git is a \textbf{Merkle DAG} --- the same hash-chained structure as blockchains and Merkle trees in distributed systems. SHA-1/SHA-256 content addressing = consistent hashing\textquotesingle s cousin for integrity.
\item
  \textbf{System Design / CI-CD:} Commits are immutable build inputs; pipelines key off SHAs. Merge queues, canary/blue-green deploys, and rollback strategies all ride on Git semantics. Feature flags connect to release engineering.
\item
  \textbf{Distributed Systems:} Git\textquotesingle s fetch/push is replication with eventual consistency between clones; divergence + merge = conflict resolution; force-push = a (dangerous) "last writer wins."
\item
  \textbf{Low-Level Design / Code Review:} Clean class design pairs with clean commit/PR design --- small, single-responsibility units in both code and history.
\item
  \textbf{Security:} Signed commits/tags = authenticity; secret scanning + history scrubbing = supply-chain hygiene; branch protection = access control.
\end{itemize}

\section{FAANG Interview Tips}\label{faang-interview-tips}

\begin{itemize}
\tightlist
\item
  \textbf{Lead with the model, not the command.} "A branch is just a pointer; commits are immutable snapshots hash-chained to their parents." This framing instantly reads as senior.
\item
  \textbf{For any \textquotesingle undo\textquotesingle{} question, state safety first:} local vs shared. Reset/rebase for local; revert/\texttt{-\/-force-with-lease} for shared. Always mention \texttt{reflog} as the safety net.
\item
  \textbf{Draw the graph.} Merge-base diagrams, before/after rebase, the three areas --- a quick ASCII sketch communicates more than paragraphs.
\item
  \textbf{Tie workflows to constraints.} Don\textquotesingle t say "Gitflow is best." Say "for continuous deploy with strong CI, trunk-based + flags; for versioned releases, Gitflow."
\item
  \textbf{Show disaster composure.} "First I\textquotesingle d \texttt{git\ reflog}, then choose the least destructive fix." Calm + correct beats fast + reckless.
\item
  \textbf{Mention guardrails unprompted.} Branch protection, required CI, \texttt{-\/-force-with-lease}, signed tags, secret scanning --- shows you think about team safety at scale.
\item
  \textbf{Use the precise vocabulary:} merge base, fast-forward, detached HEAD, dangling commit, content addressing, Merkle DAG, packfile, \texttt{ort} strategy, sparse-checkout.
\item
  \textbf{Connect to scale.} Monorepo build tooling, VFS, merge queues, stacked diffs --- signals you\textquotesingle ve thought beyond a 5-person repo.
\end{itemize}

\section{Revision Cheat Sheet}\label{revision-cheat-sheet}

\textbf{Mental model}

\begin{itemize}
\tightlist
\item
  Git = content-addressed key-value store of 4 objects: \textbf{blob} (file bytes), \textbf{tree} (dir), \textbf{commit} (tree+parent+meta), \textbf{tag} (named pointer).
\item
  Branch = 41-byte file holding a hash. HEAD = "where am I." Commits are immutable; "editing history" = new hashes + moved pointer.
\item
  Commit hash = SHA-1 of \texttt{"commit\ \textless{}len\textgreater{}\textbackslash{}0"} + text, and the text \textbf{includes parent+tree hashes} \(\rightarrow\) tamper-evident Merkle DAG.
\end{itemize}

\textbf{Three areas}

\begin{itemize}
\tightlist
\item
  Working dir \(\rightarrow\) \texttt{add} \(\rightarrow\) Index (staging) \(\rightarrow\) \texttt{commit} \(\rightarrow\) Repo. \texttt{restore}/\texttt{reset} move left. \texttt{add\ -p} for partial staging.
\end{itemize}

\textbf{Reset modes (what each touches)}

\begin{longtable}[]{@{}
  >{\raggedright\arraybackslash}p{(\columnwidth - 6\tabcolsep) * \real{0.1926}}
  >{\raggedright\arraybackslash}p{(\columnwidth - 6\tabcolsep) * \real{0.1265}}
  >{\raggedright\arraybackslash}p{(\columnwidth - 6\tabcolsep) * \real{0.1582}}
  >{\raggedright\arraybackslash}p{(\columnwidth - 6\tabcolsep) * \real{0.5227}}@{}}
\toprule\noalign{}
\rowcolor{acc!16}
\begin{minipage}[b]{\linewidth}\raggedright
Mode
\end{minipage} & \begin{minipage}[b]{\linewidth}\raggedright
HEAD
\end{minipage} & \begin{minipage}[b]{\linewidth}\raggedright
Index
\end{minipage} & \begin{minipage}[b]{\linewidth}\raggedright
Working dir
\end{minipage} \\
\midrule\noalign{}
\endhead
\bottomrule\noalign{}
\endlastfoot
\texttt{-\/-soft} & & & (changes staged) \\
\rowcolor{zebra}
\texttt{-\/-mixed} & & & (changes unstaged) \\
\texttt{-\/-hard} & & & (changes DESTROYED) \\
\end{longtable}

\textbf{Undo decision}

\begin{itemize}
\tightlist
\item
  Local, unpushed \(\rightarrow\) \texttt{reset} (soft/mixed safe; hard destroys uncommitted).
\item
  Pushed/shared \(\rightarrow\) \texttt{revert} (inverse commit; never rewrite shared history).
\item
  Lost commits \(\rightarrow\) \texttt{git\ reflog} \(\rightarrow\) \texttt{reset\ -\/-hard\ \textless{}sha\textgreater{}} / \texttt{switch\ -c\ rescue\ \textless{}sha\textgreater{}}.
\end{itemize}

\textbf{Merge vs Rebase}

\begin{itemize}
\tightlist
\item
  Merge = preserve topology, merge commit, original hashes. Rebase = linear, \textbf{new hashes}, replay. \textbf{Golden Rule: never rebase pushed/shared commits.} \texttt{pull\ -\/-rebase} for clean local history.
\end{itemize}

\textbf{Conflict}

\begin{itemize}
\tightlist
\item
  3-way: base + ours + theirs. Conflict = both changed same region. Resolve, remove markers, \texttt{add}, \texttt{-\/-continue}. \texttt{-\/-ours}/\texttt{-\/-theirs} per file. \texttt{zdiff3} + \texttt{rerere} make life easier.
\end{itemize}

\textbf{Power tools}

\begin{itemize}
\tightlist
\item
  \texttt{reflog} (recover anything), \texttt{bisect} (O(log n) regression hunt, \texttt{bisect\ run}), \texttt{cherry-pick} (backport one commit), \texttt{stash} (shelve WIP), \texttt{worktree} (parallel branches), \texttt{blame}/\texttt{log\ -S}/\texttt{log\ -L}.
\end{itemize}

\textbf{Workflows}

\begin{itemize}
\tightlist
\item
  Gitflow (versioned releases), GitHub Flow (continuous deploy), GitLab Flow (env branches), \textbf{Trunk-Based + feature flags} (FAANG scale; continuous tiny merges avoid merge hell).
\end{itemize}

\textbf{Monorepo vs Polyrepo}

\begin{itemize}
\tightlist
\item
  Mono: atomic cross-project change + uniform tooling; needs Bazel/Buck/VFS/sparse-checkout (Google/Meta). Poly: isolation + independent deploys; dependency/version coordination cost (Amazon).
\end{itemize}

\textbf{Disaster one-liners}

\begin{itemize}
\tightlist
\item
  Committed to main: \texttt{git\ branch\ feature\ \&\&\ git\ reset\ -\/-hard\ origin/main}.
\item
  Undo pushed commit: \texttt{git\ revert\ \textless{}sha\textgreater{}}.
\item
  Lost via hard reset: \texttt{git\ reflog} \(\rightarrow\) \texttt{git\ reset\ -\/-hard\ \textless{}sha\textgreater{}} (or \texttt{ORIG\_HEAD}).
\item
  Clobbered teammate: restore via reflog; prevent with \texttt{-\/-force-with-lease} + branch protection.
\item
  Leaked secret: \textbf{rotate first}, \texttt{git\ filter-repo\ -\/-path\ \textless{}f\textgreater{}\ -\/-invert-paths} / BFG, force-push, re-clone.
\item
  Detached HEAD: \texttt{git\ switch\ -c\ rescue} (before leaving) or reflog after.
\end{itemize}

\textbf{Hygiene}

\begin{itemize}
\tightlist
\item
  Atomic, single-purpose commits that compile/pass tests. Conventional Commits (\texttt{feat}/\texttt{fix}/\texttt{feat!}) \(\rightarrow\) SemVer (minor/patch/major) \(\rightarrow\) automated changelogs. Small PRs. \texttt{-\/-force-with-lease}, not \texttt{-\/-force}. Sign release tags.
\end{itemize}

\textbf{Golden sentences for the room}

\begin{enumerate}
\def\labelenumi{\arabic{enumi}.}
\tightlist
\item
  "A branch is just a pointer; Git is a Merkle DAG of immutable snapshots."
\item
  "Reset rewinds (local only); revert records an inverse (safe for shared)."
\item
  "Rebase locally for clean history; never rebase what you\textquotesingle ve pushed."
\item
  "Don\textquotesingle t panic --- \texttt{git\ reflog} first, then the least destructive fix."
\end{enumerate}

\end{document}
